\chapter{RUANG BANACH}

\section{Transformasi Alami}
 Ingat bahwa jika $V$ suatu ruang linear bernorma, maka
 \index{dual!ruang}%
 \index{ruang!dual}%
 \index{<unopurstar@$V^*$ (ruang dual dari suatu ruang linear bernorma)}%
\emph{ruang dual} $V^*$ adalah himpunan semua fungsional linear kontinu pada~$V$.  Pada $V^*$, penjumlahan dan
perkalian skalar didefinisikan secara titik demi titik. Normanya adalah norma operator biasa: jika $f \in V^*$, maka
 \[ \norm{f} := \inf\{M > 0\colon \abs{f(x)} \le M\norm{x}\quad \text{untuk setiap $x \in V$}\} = \sup\{\abs{f(x)}\colon \norm{x} \le 1\}. \]
Di bawah metrik yang diinduksi oleh norma ini, ruang dual $V^*$ lengkap dan karena itu merupakan ruang Banach.

\begin{defn}\label{defn_nls_adjoint} Jika $T\colon V \sto W$ suatu pemetaan linear kontinu antara ruang-ruang linear bernorma,
 \index{<unopurstar@$T^*$ (adjoin ruang linear bernorma)}%
$T^*\colon W^* \sto V^*$ didefinisikan seperti pada ruang vektor: $T^*(g) = g\,T$ untuk setiap
$g \in W^*$.  Pemetaan $T^*$ disebut
 \index{adjoin!pemetaan linear!antara ruang linear bernorma}%
 \index{linear!pemetaan!adjoin dari}%
\df{adjoin} dari~$T$.
\end{defn}

\begin{exam}\label{exam_dual_functor} Adjoin $T^*$ dari suatu pemetaan linear terbatas antara ruang-ruang linear bernorma $V$ dan
$W$ sendiri merupakan pemetaan linear terbatas antara $W^*$ dan~$V^*$. Dalam kategori ruang linear bernorma dan pemetaan linear
 \index{funktor!dualitas}%
 \index{funktor dualitas}%
kontinu, pasangan pemetaan $V \mapsto V^*$ (pemetaan objek) dan $T \mapsto T^*$ (pemetaan morfisme) merupakan suatu funktor
kontravarian dari $\cat{NLS_\infty}$ ke dirinya sendiri.
\end{exam}

Perhatikan bahwa Definisi~\ref{defn_nls_adjoint} berbeda dari definisi adjoin suatu pemetaan
antara ruang-ruang Hilbert (lihat~\ref{def000926}).  Karena setiap ruang Hilbert merupakan ruang linear bernorma, hal ini dapat saja
menimbulkan kebingungan.  Ada baiknya kita meyakinkan diri tentang fakta sederhana berikut: karena anti-isomorfisme (dan akibatnya
isomorfisme) antara ruang Hilbert dan dualnya yang diberikan oleh \emph{teorema Riesz--Fr\'echet} (lihat~\ref{0022057}
dan~\ref{cor_Hsp_iso_dual}), adjoin ruang linear bernorma dari suatu pemetaan linear terbatas \emph{sangat} mirip (meskipun
tentu tidak identik) dengan adjoin ruang Hilbertnya.

\begin{exam}\label{C057717}
Dual $V^{**}$ dari dual $V^*$ suatu ruang linear bernorma $V$ adalah \df{dual kedua} dari~$V$
 \index{dual!kedua}%
dan $T^{**} := \bigl(T^*\bigr)^*$ merupakan pemetaan linear terbatas dari $V^{**}$ ke~$W^{**}$.  Pada kategori
$\cat{BAN_{\boldsymbol\infty}}$ yang terdiri atas ruang Banach dan transformasi linear terbatas,
 \index{kedua!dual!funktor untuk ruang linear bernorma}%
 \index{funktor!dual kedua!untuk ruang linear bernorma}%
\df{funktor dual kedua} mencakup pemetaan objek $B \mapsto B^{**}$ dan pemetaan morfisme $T \mapsto T^{**}$.
Funktor ini kovarian.
\end{exam}

\begin{defn}\label{defn_nat_transf} Misalkan $\cat A$ dan $\cat B$ kategori, serta $F$,$G \colon \cat A \sto \cat B$
funktor kovarian.  Suatu
 \index{alami!transformasi}%
 \index{transformasi!alami}%
\df{transformasi alami} dari $F$ ke $G$ adalah pemetaan $\tau$ yang memasangkan kepada setiap objek
$A$ dalam $\cat A$ suatu morfisme $\tau_A \colon F(A) \sto G(A)$ dalam $\cat B$ sedemikian sehingga
untuk setiap morfisme $\alpha \colon A \sto A'$ dalam $\cat A$, diagram berikut komutatif.
   \[ \CD
      F(A) @>  F(\alpha)>>  F(A')  \\
         @V\tau_AVV                @VV\tau_{A'}V  \\
      G(A) @>> G(\alpha)>  G(A')
   \endCD \]
Kita menyatakan transformasi semacam itu dengan $\tau \colon F \sto G$. (Definisi transformasi alami
antara dua funktor kontravarian seharusnya jelas: cukup balikkan
panah-panah horizontal dalam diagram sebelumnya.)

Suatu transformasi alami $\tau \colon  F \sto G$ disebut
 \index{alami!ekuivalensi}%
 \index{ekuivalensi!alami}%
\df{ekuivalensi alami} jika setiap morfisme $\tau_A$ merupakan isomorfisme dalam~$\cat B$.
\end{defn}

\begin{defn}\label{defn_nat_embed} Jika $V$ suatu ruang linear bernorma, maka pemetaan $\sbsb{\tau}V \colon V \sto V^{**} \colon x
\mapsto x^{**}$, dengan $x^{**}(f) = f(x)$ untuk setiap $f \in V^*$, disebut
 \index{tauv@$\sbsb{\tau}V$ (pembenaman alami ruang linear bernorma)}%
 \index{alami!pembenaman!untuk ruang linear bernorma}%
 \index{pembenaman!alami}%
\df{pembenaman alami} dari $V$ ke~$V^{**}$. (Tidak sepenuhnya jelas bahwa pemetaan ini
injektif. Kita akan menunjukkannya dalam Proposisi~\ref{C067137}.)
\end{defn}

\begin{exam}\label{exam_2dual_nat} Dalam kategori $\cat{BAN_{\boldsymbol\infty}}$ yang terdiri atas ruang Banach
dan transformasi linear terbatas, misalkan $I$ adalah funktor identitas dan $(\,\cdot\,)^{**}$ adalah
funktor dual kedua.  Maka
 \index{alami!pembenaman!sebagai funktor}%
 \index{funktor!pembenaman alami sebagai}%
 \index{taub@$\tau_B$ (pembenaman alami $B$ ke $B^{**}$)}%
\emph{pembenaman alami} $\tau$ merupakan transformasi alami dari $I$ ke~$(\,\cdot\,)^{**}$.
\end{exam}

\begin{exam} Jelaskan secara tepat apa yang dimaksud ketika orang mengatakan bahwa isomorfisme
    \[ \fml C(X \uplus Y) \cong \fml C(X) \times \fml C(Y) \]
merupakan ekuivalensi alami.  Kategori-kategori yang terlibat adalah $\cat{CpH}$ (ruang Hausdorff kompak
dan pemetaan kontinu) serta $\cat{BAN_\infty}$ (ruang Banach dan transformasi linear
terbatas).
\end{exam}

\begin{exam}\label{C057724} Pada kategori $\cat{CpH}$ yang terdiri atas ruang Hausdorff kompak dan
pemetaan kontinu, misalkan $I$ funktor identitas dan $\fml C^*$ funktor yang
memetakan ruang Hausdorff kompak $X$ ke ruang Banach $\fml (C(X))^*$, serta memetakan suatu
fungsi kontinu $\phi\colon X \sto Y$ antara ruang-ruang Hausdorff kompak ke
transformasi linear kontraktif $\fml C^*\phi = (\fml C\phi)^*$. Definisikan
 \index{evaluasi!pemetaan}%
\emph{pemetaan evaluasi} $E$ dengan
  \[ E_X\colon X \sto \fml C^*(X)\colon x \mapsto E_X(x) \]
dengan $(E_X(x))(f) = f(x)$ untuk setiap $f \in \fml C(X)$.  (Notasi ini dapat merepotkan.  Dalam praktik, biasanya orang menulis
$E_x$ saja untuk $E_X(x)$.  Jadi kita memperoleh $E_x(f) = f(x)$, yang tentu tampak lebih baik daripada $(E_X(x))(f) = f(x)$.  Dalam kebanyakan
konteks, penyederhanaan notasi ini tampaknya tidak mungkin menimbulkan kebingungan.)
 \begin{enumerate}
  \item[(a)]  Perhatikan bahwa definisi ini memang bermakna. Jika $x \in X$, maka $E_X(x)$ benar-benar
merupakan anggota~$\fml C^*(X)$.
  \item[(b)] Tunjukkan bahwa pemetaan $E$ membuat diagram berikut komutatif:
     \[ \xymatrix@+30pt{{X}\ar[r]^*+{\phi} \ar[d]_*+{E_X}
            & {Y} \ar[d]^*+{E_Y} \\
          {\fml C^*(X)} \ar[r]_*+{\fml C^*(\phi)} & {\fml C^*(Y)}   }
     \]
 \end{enumerate}
\end{exam}

Contoh sebelumnya mungkin terasa agak mengecewakan.  Tampaknya kita sedang
berusaha membentuk suatu transformasi alami antara funktor identitas (pada
$\cat{CpH}$) dan funktor $\fml C^*$.  Masalahnya, tentu saja, ialah kedua
funktor tersebut memiliki kategori sasaran yang berbeda.  Pertanyaan ini akan kita bahas kelak.
Pemetaan evaluasi $E_X\colon X \sto \fml C^*(X)$ tentu tidak surjektif.
Perhatikan, misalnya, bahwa setiap fungsional berbentuk $E_X(x)$
mempertahankan perkalian (selain penjumlahan dan perkalian skalar). Salah satu akibatnya
ialah bahwa semua fungsional semacam itu terletak pada sfer satuan $\fml C^*(X)$.
Dengan suatu topologi baru (yang disebut \emph{topologi lemah-bintang}---lihat~\ref{C067434}) pada
$\fml C^*(X)$, jangkauan $E_X$ ternyata merupakan ruang Hausdorff kompak yang secara alami
ekuivalen dengan $X$ sendiri.  Yang lebih luar biasa lagi, pada akhirnya kita akan menunjukkan bahwa jangkauan $E_X$
dapat dipandang sebagai ruang ideal maksimal dari aljabar~$\fml C(X)$. (Lihat~\futurexref{11.2.20}{000731}.)

\begin{prop}\label{C067137} Jika $V$ suatu ruang linear bernorma, maka pembenaman alami
$\sbsb{\tau}V \colon V \sto V^{**}\colon x \mapsto x^{**}$ (yang didefinisikan dalam~\ref{defn_nat_embed})
merupakan pemetaan linear isometrik.
\end{prop}

\begin{defn} Suatu ruang Banach $B$ disebut
 \index{refleksif!ruang Banach}%
\df{refleksif} jika pembenaman alami $\sbsb{\tau}B\colon B \sto B^{**}$ (yang didefinisikan dalam proposisi sebelumnya) surjektif.
\end{defn}

Contoh~\ref{exam_2dual_nat} menunjukkan bahwa dalam kategori ruang Banach dan pemetaan linear terbatas,
pemetaan $\tau\colon B \mapsto \sbsb{\tau}B$ merupakan transformasi alami dari
funktor identitas ke funktor dual kedua.  Dengan menggunakan \emph{teorema Hahn--Banach}, kita dapat
mengatakan lebih banyak.

\begin{prop} Dalam kategori ruang Banach refleksif dan pemetaan linear terbatas, pemetaan
$\tau\colon B \mapsto \sbsb{\tau}B$ merupakan ekuivalensi alami antara funktor identitas dan
funktor dual kedua.
\end{prop}

\begin{exam} Setiap ruang Hilbert bersifat refleksif.
\end{exam}

\begin{exer} Misalkan $V$ dan $W$ ruang-ruang linear bernorma, $U$ subhimpunan konveks takkosong dari $V$,
dan $f \colon U \sto W$.  Tanpa menggunakan bentuk apa pun dari \emph{teorema nilai rata-rata}, tunjukkan
bahwa jika $df = \vc 0$ pada $U$, maka $f$ konstan pada~$U$.  (Untuk definisi $df$, yaitu
\emph{diferensial} dari~$f$, lihat Bagian 13.1--2 dalam catatan saya~\cite{Erdman:2007}.)
\end{exer}

\begin{exam}\label{C067181} Misalkan $l_\infty$ ruang Banach dari semua barisan bilangan
real yang terbatas, dan $c$ subruang $l_\infty$ yang terdiri atas semua barisan $x = (x_n)$
sedemikian sehingga $\lim_{n \sto \infty} x_n$ ada. Definisikan
  \[ f\colon c \sto \R\colon x \mapsto \lim_{n\sto \infty} x_n. \]
Maka $f$ merupakan fungsional linear kontinu pada $c$ dan $\norm f = 1$.  Berdasarkan
\emph{teorema Hahn--Banach}, $f$ memiliki perluasan ke $\wh f \in {l_\infty}^*$ dengan
$\norm{\wh f} = 1$. Untuk setiap
subhimpunan $A$ dari $\N$, definisikan barisan $\bigl(x^A_n\bigr)$ dengan $x^A_n = \left\{%
\begin{array}{ll}
    1, & \hbox{jika $n \in A$;} \\
    0, & \hbox{jika $n \notin A$.} \\
\end{array}%
\right.$ Selanjutnya, definisikan $ \mu\colon \sfml P(\N) \sto \R \colon A \mapsto \wh
f\bigl(x^A_n\bigr)$. Maka $\mu$ merupakan fungsi himpunan aditif hingga, tetapi tidak aditif
terhitung.
\end{exam}













\section{Teorema Alaoglu}
\begin{notn} Misalkan $V$ suatu ruang linear bernorma, $M \subseteq V$, dan $F \subseteq V^\ast$.
Kita definisikan
 \begin{align*}
         M^\perp &:= \{f \in V^*\colon f(x)=0 \text{ untuk setiap $x \in M$}\} \\
         F_\perp &:= \{x \in V\colon f(x)=0 \text{ untuk setiap $f \in F$}\}
  \end{align*}
$M^\perp$ dapat dibaca ``M perp'' atau ``M perp atas'', sedangkan $F_\perp$ dibaca ``F perp'' atau ``F perp bawah''.
Himpunan $M^\perp$ adalah
 \index{anihilator}%
 \index{<unopur@$M^\perp$ (anihilator suatu himpunan)}%
\df{anihilator} dari~$M$, dan $F_\perp$ adalah
 \index{praanihilator}%
 \index{<unopvr@$F_\perp$ (praanihilator suatu himpunan)}%
\df{praanihilator} dari~$F$.
\end{notn}

\begin{prop}[Anihilator]  Misalkan $B$ suatu ruang Banach, $M \subseteq B$, dan $F \subseteq B^*$. Maka
 \begin{enumerate}[label=\rm{(\alph*)}]
  \item Jika $M \subseteq N \subseteq B$, maka $N^\perp \subseteq M^\perp$.
  \item Jika $F \subseteq G \subseteq B^*$, maka $G_\perp \subseteq F_\perp$.
  \item $M^\perp$ merupakan subruang linear tertutup dari $B^*$.
  \item $F_\perp$ merupakan subruang linear tertutup dari $B$.
  \item $M \subseteq M^\perp{}_\perp$.
  \item $F \subseteq F_\perp{}^\perp$.
  \item $M^\perp{}_\perp = \bigvee M$.
  \item $M^\perp = B^*$ jika dan hanya jika $M \subseteq \{0\}$.
  \item $F_\perp = B$ jika dan hanya jika $F \subseteq \{0\}$.
  \item Jika $M$ suatu subruang linear dari $B$, maka $M^\perp = \{0\}$ jika dan hanya jika $M$
padat dalam $B$.
  \item\label{Fpp_reflexive_case} Jika $B$ refleksif, maka $F_\perp{}^\perp = \bigvee F$.
 \end{enumerate}
\end{prop}

\begin{proof}[\emph{Petunjuk pembuktian}] Untuk (k), tunjukkan bahwa $\tau(F_\perp) = F^\perp$ (dengan $\tau$
pembenaman alami) dan karena itu $F_\perp{}^\perp = F^\perp{}_\perp$.  \ns
\end{proof}

\begin{exam} Misalkan $B$ suatu ruang Banach yang \emph{tidak} refleksif.  Misalkan $F = \ran \tau$ (dengan $\tau$
pembenaman alami).  Maka $F_\perp{}^\perp = B^{**} \ne \bigvee F$, sehingga jelas bahwa~\ref{Fpp_reflexive_case}
dalam proposisi sebelumnya tidak harus berlaku dalam kasus nonrefleksif.
\end{exam}

Hasil berikut adalah analog ruang Banach dari \emph{teorema dasar aljabar linear}~\ref{00016025} dan
Proposisi~\ref{prop_ker_adj_perp_ran_op}.

\begin{prop}\label{prop_FTLA_Bsp} Misalkan $T \in \ofml B(B,C)$, dengan $B$ dan $C$ ruang-ruang Banach. Maka
 \begin{enumerate}[label=\rm{(\alph*)}]
  \item $\ker T^\ast = (\ran T)^\perp$.
  \item $\ker T = (\ran T^*)_\perp$.
\vskip 3 pt
  \item $\clo{\ran T} = (\ker T^*)_\perp$.
\vskip 3 pt
  \item $\clo{\ran T^*} \subseteq (\ker T)^\perp$.
 \end{enumerate}
\end{prop}

\begin{exam} Suatu contoh operator ruang Banach $T$ yang membuat $\clo{\ran T^*}$ dan $(\ker T)^\perp$ berbeda
diberikan dalam Latihan 7, halaman 243 dari~\cite{BrownP:1970}.
\end{exam}

Contoh berikut tentu merupakan akibat langsung yang sepele dari Proposisi~\ref{00015032}, tetapi cobalah memberikan
pembuktian langsung yang sederhana.

\begin{exam}\label{X_l2ball_not_cpt} Bola satuan tertutup dalam $l_2$ tidak kompak.
\end{exam}

Topologi norma biasa pada ruang linear bernorma ternyata terlalu besar untuk beberapa penerapan.  Seperti telah kita lihat,
topologi ini memiliki begitu banyak himpunan terbuka sehingga bola satuan tertutup tidak kompak (kecuali dalam kasus berdimensi hingga). Di
sisi lain, terlalu sedikit himpunan terbuka (sebagai contoh ekstrem, topologi indiskret) akan merusak teori dualitas karena tidak
menyediakan cukup banyak fungsional linear kontinu.  \emph{Topologi lemah-bintang} (yang didefinisikan
di bawah) pada dual suatu ruang linear bernorma ternyata tepat ukurannya.  Topologi ini memiliki cukup banyak himpunan terbuka untuk menjamin
teori dualitas yang kuat dan cukup sedikit himpunan terbuka untuk membuat bola satuan tertutup kompak
(lihat \emph{teorema Alaoglu}~\ref{C067441}).

\begin{defn}\label{C067434} Misalkan $V$ suatu ruang linear bernorma. 
 \index{lemah!topologi!pada ruang bernorma}%
 \index{topologi!lemah!pada ruang linear bernorma}%
\df{topologi lemah} pada $V$ adalah topologi lemah yang diinduksi oleh unsur-unsur~$V^*$. (Jadi,
topologi lemah pada ruang dual $V^*$ adalah topologi lemah yang diinduksi oleh
unsur-unsur~$V^{**}$.) Ketika suatu jaring $(x_\lambda)$ konvergen secara lemah ke vektor $a$ dalam $V$,
kita menulis
 \index{<arrowright@$x_\lambda \to^w a$ (konvergensi lemah)}%
$x_\lambda \to^w a$. 
 \index{w*@$w^*$-topologi (topologi lemah-bintang)}%
 \index{topologi!lemah-bintang}%
\df{topologi-$w^*$} (dibaca \emph{topologi lemah-bintang}) pada ruang dual $V^*$ adalah
topologi lemah yang diinduksi oleh unsur-unsur $\ran \tau$, dengan $\tau$ pembenaman alami
dari $V$ ke~$V^{**}$.  Ketika suatu jaring $(f_\lambda)$ konvergen secara lemah-bintang ke vektor $g$
dalam $V^*$, kita menulis
 \index{<arrowright@$f_\lambda \to^{w^*} g$ (konvergensi lemah-bintang)}%
$f_\lambda \to^{w^*} g$. Perhatikan bahwa ketika suatu ruang Banach $B$ refleksif (khususnya, untuk ruang
Hilbert), topologi lemah dan lemah-bintang pada $B^*$ identik.
\end{defn}

\begin{prop} Misalkan $f \in V^*$, dengan $V$ suatu ruang Banach.  Untuk setiap $\epsilon > 0$ dan setiap
subhimpunan hingga $F \subseteq V$, definisikan
  \[ U(f;F;\epsilon)
           = \{g \in V^*\colon \abs{f(x) - g(x)} < \epsilon \text{ untuk setiap $x \in F$}\}. \]
Keluarga semua himpunan semacam itu merupakan basis bagi topologi-$w^*$ pada~$V^*$.
\end{prop}

Hasil berikut, \emph{teorema Alaoglu}, menyatakan kekompakan-$w^*$ dari bola satuan tertutup dalam
dual suatu ruang linear bernorma $V$ dan merupakan salah satu teorema yang paling berguna dalam analisis fungsional.  Karena itu,
pada awalnya mungkin agak mengejutkan bahwa pembuktian lengkap yang dijelaskan secara cermat
begitu sulit ditemukan.  Pembuktian standar pada dasarnya sangat sederhana; intinya hanyalah mengenali
bahwa anggota-anggota suatu keluarga $\fml F$ dari fungsi bernilai skalar dapat dipandang
  \begin{enumerate}
     \item[(i)] sebagai fungsi yang didefinisikan pada hasil kali takterhitung dari subhimpunan kompak medan skalar; atau
     \item[(ii)] sebagai anggota bola satuan tertutup dari~$V^*$.
  \end{enumerate}
Inti pembuktiannya adalah meyakinkan diri bahwa konvergensi suatu jaring fungsi semacam itu dalam ruang hasil kali
tepat ``sama'' dengan konvergensi-$w^*$ jaring tersebut dalam bola satuan tertutup dari~$V^*$.  Upaya untuk sekaligus
mengidentifikasi dan membedakan dua himpunan fungsi ini menantang secara notasi.  Jadi, saran terbaik saya ialah mengerjakannya
sendiri sampai Anda melihat betapa cerdik, tetapi pada dasarnya sederhana, argumen tersebut.

\begin{thm}[Teorema Alaoglu]\label{C067441} Bola satuan tertutup dalam dual suatu ruang linear bernorma
 \index{teorema Alaoglu}%
kompak dalam topologi-$w^*$.
\end{thm}

\begin{proof}[\emph{Petunjuk pembuktian}] Misalkan $V$ suatu ruang linear bernorma, dan nyatakan dengan $[V^*]_1$ bola satuan tertutup dari
ruang dualnya.  Untuk setiap $x \in V$, definisikan $D_x = \{\lambda \in \K \colon \abs\lambda \le \norm x\}$.  Hasil kali
\smash[b]{$\mathbf P = \prod\limits_{x \in V} D_x$} kompak berdasarkan \emph{teorema Tychonoff}.  Tinjau fungsi
   \[ \Phi\colon [V^*]_1 \sto \mathbf P \colon f \mapsto \bigl(f(x)\bigr)_{x \in V} \]
dengan $[V^*]_1$ dilengkapi topologi-$w^*$ dan $\mathbf P$ dilengkapi topologi hasil kali (lihat~\ref{C021437}).  Di sini kita menulis
$\bigl(f(x)\bigr)_{x \in V}$ untuk keluarga bilangan terindeks (atau barisan tergeneralisasi), bukan notasi yang lebih lazim
$\bigl(f_x\bigr)_{x \in V}$.  (Untuk pembahasan lebih lanjut tentang keluarga terindeks, lihat~\cite{Erdman:2007}, Bagian~7.2.)  Mudah dilihat bahwa
$\Phi$ injektif.  Untuk menunjukkan bahwa pemetaan itu kontinu, misalkan $\bigl(f_\lambda\bigr)_{\lambda \in \Lambda}$ suatu jaring dalam $[V^*]_1$
dan $g$ suatu fungsi dalam $[V^*]_1$ sedemikian sehingga $f_\lambda \to^{w^*} g$.  Perhatikan bahwa hal ini terjadi jika dan hanya jika
$\pi_x(f_\lambda) \sto \pi_x(g)$ untuk setiap $x \in V$ (dengan $\pi_x$ proyeksi koordinat pada~$\mathbf P$---lihat~\ref{C015531}).
Terakhir, misalkan $\bigl(f_\lambda\bigr)_{\lambda \in \Lambda}$ suatu jaring dalam $[V^*]_1$ sedemikian sehingga $\bigl(\Phi(f_\lambda)\bigr)_{\lambda \in \Lambda}$
konvergen dalam~$\mathbf P$.  Tunjukkan bahwa jaring $\bigl(f_\lambda(x)\bigr)_{\lambda \in \Lambda}$ konvergen dalam $\K$ untuk setiap $x \in V$.
Misalkan $g(x) = \lim\limits_\lambda f_\lambda(x)$ untuk setiap $x \in V$, lalu buktikan bahwa $g$ linear, $\norm g \le 1$,
$f_\lambda \to^{w^*} g$ dalam $[V^*]_1$, dan $\Phi(f_\lambda) \sto \Phi(g)$ dalam~$\mathbf P$.  \ns
\end{proof}

Sepengetahuan saya, versi pembuktian standar ini yang paling ramah pembaca terdapat dalam~\cite{BrownP:1977}, Teorema 15.11.
Pembuktian yang jauh lebih singkat, jauh lebih mudah secara notasi, jauh lebih elegan, dan jauh kurang mencerahkan tersedia dalam~\cite{Pedersen:1995}, Teorema 2.5.2.
Menurut saya, pembuktian itu kurang mencerahkan karena bergantung pada jaring universal.  Suatu jaring dalam himpunan $S$ disebut
 \index{universal!jaring}%
 \index{jaring!universal}%
\df{universal} jika untuk setiap subhimpunan $T$ dari $S$, jaring tersebut pada akhirnya berada dalam $T$ atau pada akhirnya berada dalam komplemennya. Meskipun
setiap jaring dapat dibuktikan memiliki subjaring universal, bahkan mencari satu contoh nontrivial dari jaring universal (yakni, jaring yang tidak
pada akhirnya konstan) memerlukan \emph{aksioma pilihan}.






\section{Teorema Pemetaan Terbuka}\label{C0694}
Salah satu pertanyaan yang dapat (dan patut) diajukan tentang setiap kategori konkret ialah apakah setiap morfisme bijektif selalu merupakan
isomorfisme.  Dalam beberapa kategori jawabannya jelas \,\emph{ya} (misalnya, dalam $\cat{VEC}$).  Dalam kategori lain
jawabannya jelas \,\emph{tidak} (dalam $\cat{TOP}$, misalnya, tinjau pemetaan identitas dari bilangan
real bertopologi diskret ke bilangan real bertopologi biasa).  Dalam kategori-kategori yang muncul dalam
kajian analisis, aljabar sering menarik dengan penuh semangat ke satu arah, sementara topologi melawan dengan kuat dan menarik
ke arah yang lain.  Di sini pertanyaan tersebut dapat menjadi sangat menarik---dan jawabannya sama sekali tidak sepele.  Dalam kategori
$\cat{BAN_\infty}$ yang terdiri atas ruang Banach dan pemetaan linear kontinu, pertanyaan ini ternyata memiliki kedalaman yang menarik.  Dalam
kasus ini aljabar menang; jawabannya ya.  Hal ini ternyata merupakan akibat langsung dari \emph{teorema
pemetaan terbuka} yang termasyhur.  Ketika Anda menelaah pembuktian hasil berikut, perhatikan bahwa untuk pemetaan yang dibahas,
kelengkapan domain dan kelengkapan kodomain sama-sama esensial---dan dengan alasan yang sangat berbeda.

\begin{defn} Suatu pemetaan $f \colon X \sto Y$ antara ruang-ruang topologis disebut
 \index{terbuka!pemetaan}%
 \index{pemetaan!terbuka}%
\df{terbuka} jika memetakan himpunan terbuka ke himpunan terbuka; yaitu, $f$ merupakan pemetaan terbuka apabila $f^{\sto}(U)$ terbuka dalam $Y$
setiap kali $U$ terbuka dalam~$X$.
\end{defn}

\emph{Teorema pemetaan terbuka} menyatakan bahwa surjeksi linear kontinu antara ruang-ruang Banach selalu merupakan pemetaan terbuka.
Ada satu rincian teknis yang membuat pembuktian hasil ini agak rumit.  Misalkan $T$ suatu pemetaan linear kontinu
dari ruang Banach $B$ ke ruang linear bernorma~$V$, dan $R$ adalah citra oleh $T$ dari suatu bola terbuka di sekitar
titik asal dalam~$B$.  Kita perlu mengetahui bahwa jika tutupan $R$ memuat suatu bola terbuka di sekitar titik asal dalam $V$, maka
$R$ sendiri memuat bola yang sama. Akan memudahkan jika bagian informasi yang krusial ini kita pisahkan sebagai proposisi berikut.

\begin{notn} Dalam ruang linear bernorma $V$, kita menyatakan dengan
 \index{<unop@$(V)_r$ (bola terbuka berjari-jari $r$ di sekitar titik asal)}%
$(V)_r$ bola terbuka di sekitar titik asal dalam $V$ yang berjari-jari~$r$.
\end{notn}

\begin{prop}\label{prop_for_OMT} Misalkan $T\colon B \sto V$ suatu pemetaan linear kontinu dari ruang Banach ke ruang linear bernorma.
Jika $r$, $s > 0$ dan $(V)_s \subseteq \clo{T\bigl(\,(B)_r\,\bigr)}$, maka $(V)_s \subseteq T\bigl(\,(B)_r\,\bigr)$.
\end{prop}

\begin{proof}[\emph{Petunjuk pembuktian}] Pertama, buktikan bahwa tanpa mengurangi keumuman kita dapat mengasumsikan $r = s = 1$.  Dengan demikian,
hipotesis kita adalah $(V)_1 \subseteq \clo{T\bigl(\,(B)_1\,\bigr)}$.

Misalkan $z \in (V)_1$.  Pembuktian selesai jika kita dapat menunjukkan bahwa $z \in T\bigl(\,(B)_1\,\bigr)$.  Pilih $\delta > 0$ sedemikian
sehingga $\norm z < 1 - \delta$, dan misalkan $y = (1 - \delta)^{-1}z$.  Bangun suatu barisan $(v_k)_{k=0}^\infty$ dari vektor-vektor dalam~$V$
yang memenuhi
   \begin{enumerate}
     \item[(i)] $\norm{y - v_k} < \delta^k$ untuk $k = 0,1,2,\dots$; dan
     \item[(ii)] $v_k - v_{k-1} \in T\bigl(\,\delta^{k-1}(B)_1\,\bigr)$ untuk $k = 1,2,3,\dots$.
   \end{enumerate}
Untuk melakukannya, mulailah dengan $v_0 = \vc 0$. Gunakan hipotesis untuk menyimpulkan bahwa $y - v_0 \in \clo{T\bigl(\,(B)_1\,\bigr)}$ dan
karena itu terdapat vektor $w_0 \in T\bigl(\,(B)_1\,\bigr)$ sedemikian sehingga $\norm{(y - v_0) - w_0} < \delta$.  Misalkan
$v_1 = v_0 + w_0$, lalu periksa bahwa (i) dan (ii) berlaku untuk $k = 1$.

Kemudian cari $w_1 \in T\bigl(\,\delta (B)_1\,\bigr)$ sedemikian sehingga $\norm{(y - v_1) - w_1} < \delta^2$, dan misalkan $v_2 = v_1 + w_1$.
Periksa bahwa (i) dan (ii) berlaku untuk $k= 2$.  Lanjutkan secara induktif.

Untuk setiap $n \in \N$, pilih vektor $u_n \in \delta^{n-1}(B)_1$ sedemikian sehingga $T(u_n) = v_n - v_{n-1}$. Tunjukkan bahwa barisan
$(u_n)$ dapat dijumlahkan (gunakan~\ref{prop_abs_sum_implies_sum}), dan misalkan $x = \sum_{n=1}^\infty u_n$.  Selesaikan pembuktian dengan menunjukkan
bahwa $Tx = y$ dan bahwa $z \in T\bigl(\,(B)_1\,\bigr)$.   \ns
\end{proof}

\begin{thm}[Teorema pemetaan terbuka]\label{C069414} Setiap surjeksi linear terbatas antara
 \index{terbuka!pemetaan!teorema}%
 \index{morfisme bijektif!dapat diinverskan!dalam $\cat{BAN_\infty}$}%
ruang-ruang Banach merupakan pemetaan terbuka.
\end{thm}

\begin{proof}[\emph{Petunjuk pembuktian}] Misalkan $T\colon B \sto C$ suatu surjeksi linear terbatas antara ruang-ruang Banach.
Mulailah dengan membuktikan klaim berikut: yang perlu kita tunjukkan hanyalah bahwa $T\bigl((B)_1\bigr)$ memuat suatu bola terbuka di sekitar
titik asal dalam~$C$.  Untuk sementara, misalkan kita mengetahui bahwa bola semacam itu ada.  Namai bola itu~$W$.  Diberikan suatu subhimpunan terbuka $U$ dari $B$, kita
ingin menunjukkan bahwa citranya oleh $T$ terbuka dalam~$C$.  Untuk itu, ambil sebarang titik $y$ dalam $T(U)$ dan pilih $x$
dalam $U$ sedemikian sehingga $y = Tx$.  Karena $U - x$ merupakan lingkungan titik asal dalam $B$, kita dapat mencari bilangan $r > 0$ sedemikian sehingga
$r\,(B)_1 \subseteq U - x$.  Buktikan bahwa $T(U)$ terbuka dengan menunjukkan bahwa $y + rW \subseteq T(U)$.

Untuk membuktikan klaim di atas, misalkan $V = T\bigl((B)_1\bigr)$ dan perhatikan bahwa karena $T$ surjektif, $C = T(B) =
\bigcup_{k=1}^\infty k\clo V$.  Gunakan versi \emph{teorema kategori Baire} yang menyatakan bahwa jika suatu ruang metrik lengkap takkosong
merupakan gabungan suatu barisan himpunan tertutup, maka sekurang-kurangnya salah satu himpunan tersebut memiliki interior takkosong.
(Lihat, misalnya, Teorema 29.3.21 dalam~\cite{Erdman:2007}.)  Pilih $k \in \N$ sedemikian sehingga $\bigl(k\clo V\bigr)^o \neq \emptyset$.
Dengan demikian, terdapat $y \in \clo V$ dan $r > 0$ sedemikian sehingga $y + (C)_r \subseteq k\clo V$.  Dari kekonveksan
$k\clo V$, tidak sulit melihat bahwa $(C)_r \subseteq k\clo V$.  (Tuliskan sebarang titik $z$ dalam $(C)_r$ sebagai
$\frac12(y + z) + \frac12(-y + z)$.) Sekarang diperoleh $(C)_{r/k} \subseteq \clo V$.  Gunakan Proposisi~\ref{prop_for_OMT}.
\ns \end{proof}

\begin{cor}\label{C069417} Setiap bijeksi linear terbatas antara ruang-ruang Banach merupakan suatu
isomorfisme.
\end{cor}

\begin{cor} Setiap surjeksi linear terbatas antara ruang-ruang Banach
 \index{BAN@$\cat{BAN_\infty}$!hasil bagi dalam}%
 \index{hasil bagi!pemetaan!dalam $\cat{BAN_\infty}$}%
merupakan pemetaan hasil bagi dalam~$\cat{BAN_\infty}$.
\end{cor}

\begin{exam} Misalkan $\fml C_u$ himpunan fungsi kontinu bernilai real pada $[0,1]$ dengan
norma seragam, dan $\fml C_1$ himpunan fungsi yang sama dengan norma $\mathrm{L_1}$
(\,$\norm{f}_1 = \int_0^1\abs{f}\,d\lambda$\,). Maka pemetaan identitas $I\colon \fml C_u
\sto \fml C_1$ merupakan bijeksi linear kontinu, tetapi tidak terbuka. (Mengapa hal ini tidak
bertentangan dengan \emph{teorema pemetaan terbuka}?)
\end{exam}

\begin{exam} Misalkan $l$ himpunan barisan bilangan real yang pada akhirnya bernilai nol,
$V$ adalah $l$ yang dilengkapi norma $l_1$, dan $W$ adalah $l$ dengan norma seragam.  Pemetaan
identitas $I\colon V \sto W$ merupakan bijeksi linear terbatas, tetapi bukan suatu isomorfisme.
\end{exam}

\begin{exam} Pemetaan $T\colon \R^2 \sto \R\colon (x,y) \mapsto x$ menunjukkan bahwa meskipun
surjeksi linear terbatas antara ruang-ruang Banach memetakan himpunan terbuka ke himpunan terbuka, pemetaan tersebut tidak harus
memetakan himpunan tertutup ke himpunan tertutup.
\end{exam}

\begin{prop} Misalkan $V$ suatu ruang vektor, dan misalkan $\norm{\cdot}_1$ dan
$\norm{\cdot}_2$ adalah norma-norma pada~$V$ dengan topologi masing-masing $\ofml T_1$
dan~$\ofml T_2$.  Jika $V$ lengkap terhadap kedua norma dan jika $\ofml T_1
\supseteq \ofml T_2$, maka $\ofml T_1 = \ofml T_2$.
\end{prop}

\begin{exer} Apakah terdapat suatu barisan $(a_n)$ dari bilangan kompleks yang memenuhi
syarat berikut: suatu barisan $(x_n)$ dalam $\C$ dapat dijumlahkan secara mutlak jika dan hanya jika
$(a_nx_n)$ terbatas? \emph{Petunjuk.} Tinjau $ T\colon l_\infty \sto l_1\colon (x_n)
\mapsto (x_n/a_n)$.
\end{exer}

\begin{exam}\label{C069431} Misalkan $\fml C^2([a,b])$ keluarga semua fungsi bernilai real yang dua kali
dapat didiferensialkan secara kontinu pada interval $[a,b]$. Pandang keluarga ini sebagai ruang vektor
dengan cara biasa. Untuk setiap $f \in \fml C^2([a,b])$, definisikan
    \[ \norm{f} = \norm{f}_u + \norm{f'}_u + \norm{f''}_u\,. \]
Ini memang suatu norma, dan dengan norma ini $\fml C^2([a,b])$ merupakan ruang Banach.
\end{exam}

\begin{exam}\label{C069432} Misalkan $\fml C^2([a,b])$ ruang Banach yang diberikan dalam
Contoh~\ref{C069431}, dan misalkan $p_0$ dan $p_1$ merupakan anggota $\fml C([a,b])$.   Definisikan
  \[T\colon C^2([a,b]) \sto C([a,b])\colon
                      f \mapsto f'' + p_1f' + p_0f.\]
Maka $T$ merupakan pemetaan linear terbatas.
\end{exam}

\begin{exer} Tinjau suatu persamaan diferensial berbentuk
   \[ y'' + p_1\,y' + p_0\,y = q  \tag{$\ast$} \]
dengan $p_0$, $p_1$, dan $q$ fungsi-fungsi bernilai real kontinu (yang tetap) pada interval
$[a,b]$.  Nyatakan secara tepat, lalu buktikan, pernyataan bahwa \emph{solusi-solusi $(\ast)$
bergantung secara kontinu pada nilai awal}.  Anda boleh menggunakan tanpa bukti teorema standar berikut:
untuk setiap titik $c$ dalam $[a,b]$ dan setiap pasangan bilangan real $a_0$ dan $a_1$,
terdapat solusi tunggal untuk $(\ast)$ sedemikian sehingga $y(c) = a_0$ dan $y'(c) = a_1$.
\emph{Petunjuk.} Untuk $c \in [a,b]$ yang tetap, tinjau pemetaan
  \[ S\colon \fml C^2([a,b]) \sto \fml C([a,b]) \times \R^2 \colon
                                        f \mapsto \bigl(Tf,f(c),f'(c)\bigr) \]
dengan $\fml C^2([a,b])$ ruang Banach dari Contoh~\ref{C069431} dan $T$ pemetaan
linear terbatas dari Contoh~\ref{C069432}.
\end{exer}

\begin{defn} Suatu pemetaan linear terbatas $T\colon V \sto W$ antara ruang-ruang linear bernorma disebut
 \index{terbatas!jauh dari nol}%
\df{terbatas jauh dari nol} (atau
 \index{terbatas!dari bawah}%
\df{terbatas dari bawah}) jika terdapat bilangan $\delta > 0$ sedemikian sehingga $\norm{Tx} \ge
\delta\norm{x}$ untuk setiap $x \in V$. (Ingat bahwa kata ``terbatas'' memiliki satu arti ketika
menerangkan suatu pemetaan linear dan arti yang sangat berbeda ketika diterapkan pada fungsi umum;
demikian pula ungkapan ``terbatas jauh dari nol''. Lihat Definisi~\ref{def_bdd_away_zero}.)
\end{defn}

\begin{prop} Suatu pemetaan linear terbatas antara ruang-ruang Banach bersifat terbatas jauh dari nol jika dan
hanya jika pemetaan itu memiliki jangkauan tertutup dan kernel nol.
\end{prop}


















\section{Teorema Graf Tertutup}
Syarat perlu dan cukup yang sederhana agar suatu pemetaan linear $T\colon B \sto C$ antara ruang-ruang Banach
bersifat kontinu ialah grafnya tertutup.  Dalam satu arah, hal ini hampir jelas.  Kebalikannya
disebut \emph{teorema graf tertutup}.  Ketika kita mensyaratkan graf $T$
 \index{graf!tertutup}%
``tertutup'', maksud kita ialah tertutup sebagai subhimpunan $B \times C$ dengan topologi hasil kali.  Ingat (dari~\ref{exam_prod_top_norm})
bahwa topologi pada jumlah langsung $B \oplus C$ yang diinduksi oleh norma hasil kali sama dengan topologi hasil kali pada $B \times C$.

\begin{exer} Jika $T\colon B \sto C$ suatu pemetaan linear kontinu antara ruang-ruang Banach, maka grafnya tertutup.
\end{exer}

\begin{thm}[Teorema graf tertutup]\label{thm_CGT} Misalkan $T\colon B \sto C$ suatu pemetaan linear antara ruang-ruang Banach.
 \index{tertutup!teorema graf}%
Jika graf $T$ tertutup, maka $T$ kontinu.
\end{thm}

\begin{proof}[\emph{Petunjuk pembuktian}] Misalkan $G = \{(x,Tx)\colon x \in B\}$ graf~$T$. Terapkan
(Korolar~\ref{C069417} dari) \emph{teorema pemetaan terbuka} pada pemetaan
  \[ \pi\colon G \sto B\colon (x,Tx) \mapsto x\,. \]
\ns \end{proof}

Proposisi berikut memberikan syarat perlu dan cukup yang sangat sederhana agar graf suatu pemetaan linear antara ruang-ruang Banach
tertutup.

\begin{prop}\label{prop_nasc_closed_graph} Misalkan $T\colon B \sto C$ suatu pemetaan linear antara ruang-ruang Banach. Maka $T$
kontinu jika dan hanya jika syarat berikut dipenuhi:
   \[ \text{jika $x_n \sto 0$ dalam $B$ dan $Tx_n \sto c$ dalam $C$, maka $c = 0$.} \]
\end{prop}

\begin{prop} Misalkan $T \colon B \sto C$ suatu pemetaan linear antara ruang-ruang Banach sedemikian sehingga jika
$x_n \sto 0$ dalam $B$, maka $(g \circ T)(x_n) \sto 0$ untuk setiap $g \in C^*$.  Maka $T$
terbatas.
\end{prop}

\begin{prop} Misalkan $T \colon B \sto C$ suatu fungsi linear antara ruang-ruang Banach. Definisikan
$T^*f = f \circ T$ untuk setiap $f \in C^*$.  Jika $T^*$ memetakan $C^*$ ke dalam $B^*$, maka $T$
kontinu.
\end{prop}

\begin{defn}\label{C069624} Suatu pemetaan linear $T$ dari ruang vektor ke ruang itu sendiri disebut
 \index{idempoten}%
\df{idempoten} jika $T^2 = T$.
\end{defn}

\begin{prop}\label{prop_cont_idem_op} Suatu pemetaan linear idempoten dari ruang Banach ke ruang itu sendiri
kontinu jika dan hanya jika kernel dan jangkauannya tertutup.
\end{prop}

\begin{proof}[\emph{Petunjuk pembuktian}] Untuk bagian ``jika'', gunakan Proposisi~\ref{prop_nasc_closed_graph}.  \ns
\end{proof}

Proposisi berikut merupakan penerapan \emph{teorema graf tertutup} yang berguna pada operator ruang Hilbert.

\begin{prop} Misalkan $H$ suatu ruang Hilbert, serta $S$ dan $T$ fungsi dari $H$ ke dirinya sendiri sedemikian sehingga
    \[ \langle Sx,y \rangle = \langle x,Ty \rangle \]
untuk setiap $x$, $y \in H$.  Maka $S$ dan $T$ merupakan operator linear terbatas.
\end{prop}

\begin{proof}[\emph{Petunjuk pembuktian}] Untuk membuktikan linearitas pemetaan $S$, tinjau hasil kali dalam dari vektor-vektor
$S(\alpha x + \beta y) - \alpha Sx - \beta Sy$ dan $z$, dengan $x$, $y$, $z \in H$ dan $\alpha$, $\beta \in \K$.
Untuk membuktikan kekontinuan $S$, tunjukkan bahwa grafnya tertutup.  \ns
\end{proof}

Ingat bahwa dalam proposisi sebelumnya, operator $T$ adalah \emph{adjoin (ruang Hilbert)} dari $S$;
yaitu, $T = S^*$.

\begin{defn} Misalkan $a < b$ dalam~$\R$. Suatu fungsi $f\colon [a,b] \sto \R$ disebut
 \index{dapat didiferensialkan secara kontinu}%
 \index{diferensiabel!secara kontinu}%
dapat didiferensialkan secara kontinu jika fungsi itu dapat didiferensialkan pada (suatu himpunan terbuka yang memuat) $[a,b]$
dan turunannya $f\,'$ kontinu pada~$[a,b]$.  Himpunan semua fungsi bernilai real yang dapat
didiferensialkan secara kontinu pada $[a,b]$ dinotasikan dengan $\fml C^1(\,[a,b]\,)$.
\end{defn}

\begin{exam} Misalkan $D\colon \fml C^1(\,[0,1]\,) \sto C(\,[0,1]\,)\colon f \mapsto f\,'$, dengan
baik $C^1(\,[0,1]\,)$ maupun $C(\,[0,1]\,)$ dilengkapi norma seragam. Pemetaan
$D$ linear dan memiliki graf tertutup, tetapi tidak kontinu.  (Mengapa hal ini tidak bertentangan dengan
\emph{teorema graf tertutup}?  Apa yang dapat kita simpulkan tentang~$\fml C^1(\,[0,1]\,)$?)
\end{exam}

\begin{prop} Misalkan $S\colon A \sto B$ dan $T\colon B \sto C$ pemetaan linear antara
ruang-ruang Banach, serta $T$ terbatas dan injektif.  Maka $S$ terbatas jika dan hanya jika $TS$ terbatas.
\end{prop}











\section{Dualitas Ruang Banach}

\begin{conv}\label{conv_subsp_Bsp} Dalam konteks ruang Banach, kita menggunakan konvensi yang sama untuk pemakaian kata
``subruang'' seperti yang kita gunakan untuk ruang Hilbert. Kata ``subruang'' akan selalu berarti
 \index{subruang!dari ruang Banach}%
 \index{konvensi!subruang ruang Banach bersifat tertutup}%
\emph{subruang vektor tertutup}. Untuk menunjukkan bahwa $M$ adalah subruang dari suatu ruang Banach $B$, kita menulis
 \index{<binrelle@$\preccurlyeq$ (subruang dari ruang Banach atau Hilbert)}%
$M \preccurlyeq B$.
\end{conv}

\begin{prop} Jika $M$ adalah subruang dari suatu ruang Banach $B$,
 \index{hasil bagi!ruang Banach}%
 \index{Banach!ruang!hasil bagi}%
 \index{ruang!hasil bagi!Banach}%
maka ruang linear bernorma hasil bagi $B/M$ juga merupakan ruang Banach.
\end{prop}

\begin{proof}[\emph{Petunjuk pembuktian}] Untuk menunjukkan bahwa $B/M$ lengkap, cukup ditunjukkan bahwa dari
sebarang barisan Cauchy $(u_n)$ dalam $B/M$ kita dapat mengambil suatu subbarisan yang konvergen. Untuk itu,
perhatikan bahwa jika $(u_n)$ adalah barisan Cauchy, maka untuk setiap $k \in \N$ terdapat $N_k \in \N$ sedemikian sehingga
$\norm{u_n - u_m} < 2^{-k}$ setiap kali $m$, $n \ge N_k$. Sekarang gunakan induksi. Pilih $n_1 = N_1$.
Setelah $n_1$, $n_2$, \dots, $n_k$ dipilih sedemikian sehingga $n_1 < n_2 < \dots < n_k$ dan $N_j \le n_j$ untuk
$1 \le j \le k$, pilih $n_{k+1} = \max\{N_{k+1}, n_k + 1\}$. Dengan demikian diperoleh subbarisan
$\bigl(u_{n_k}\bigr)$ yang memenuhi $\norm{u_{n_k} - u_{n_{k+1}}} < 2^{-k}$ untuk setiap $k \in \N$.
Dari setiap kelas ekuivalensi $u_{n_k}$, pilih suatu vektor $x_k$ dalam $B$ sedemikian sehingga $\norm{x_k - x_{k+1}} < 2^{-k+1}$.
Kemudian Proposisi~\ref{prop_abs_sum_implies_sum} menunjukkan bahwa barisan $(x_k - x_{k+1})$ dapat
dijumlahkan. Dari sini, simpulkan bahwa barisan $(x_k)$ dan $\bigl(u_{n_k}\bigr)$ konvergen.
(Rincian argumen ini dituliskan dengan baik dalam pembuktian Teorema 3.4.13 pada~\cite{BrownP:1970}.) \ns
\end{proof}

\begin{prop} Jika $J$ adalah ideal tertutup dalam suatu aljabar Banach $A$,
 \index{hasil bagi!aljabar Banach}%
 \index{Banach!aljabar!hasil bagi}%
 \index{aljabar!hasil bagi!Banach}%
maka $A/J$ merupakan aljabar Banach.
\end{prop}

\begin{thm}[Teorema hasil bagi fundamental untuk $\cat{BALG}$] Misalkan $A$ dan $B$ aljabar Banach
 \index{BALG@$\cat{BALG}$!hasil bagi dalam}%
 \index{hasil bagi!teorema!untuk $\cat{BALG}$}%
dan $J$ ideal tertutup sejati dalam~$A$. Jika $\phi$ suatu homomorfisme dari $A$ ke $B$ dan
$\ker\phi \supseteq J$, maka terdapat homomorfisme tunggal $\wt\phi\colon A/J \sto
B$ yang membuat diagram berikut komutatif.
   \[ \xymatrix@+30pt{A \ar[d]_*+{\pi} \ar[dr]^*+{\phi} & \\
                           A/J \ar[r]_*+{\widetilde\phi} & B  }\]
Selain itu, $\wt\phi$ injektif jika dan hanya jika $\ker\phi = J$; dan $\wt\phi$
surjektif jika dan hanya jika $\phi$ surjektif.
\end{thm}

\begin{prop}\label{C037821} Misalkan $J$ ideal tertutup sejati dalam suatu aljabar Banach komutatif
beridentitas~$A$. Maka $J$ maksimal jika dan hanya jika $A/J$ adalah medan.
\end{prop}

Proposisi berikut menyatakan syarat-syarat yang cukup sederhana agar suatu epimorfisme antara ruang Banach
ternyata menjadi faktor dalam morfisme-morfisme lain. Sepintas lalu, kegunaan hasil semacam ini
mungkin tampak agak meragukan. Namun, ketika membuktikan Proposisi~\ref{00151231}, perhatikan bahwa
inilah \emph{tepatnya} yang Anda perlukan.

\begin{prop}\label{prop_surj_as_factor} Tinjau kategori ruang Banach
dan pemetaan linear kontinu. Jika $T\colon A \sto C$, $R\colon A \sto B$, $R$
surjektif, dan $\ker R \subseteq \ker T$, maka terdapat
$S\colon B \sto  C$ tunggal sedemikian sehingga $SR = T$. Selain itu, $S$ injektif jika dan
hanya jika $\ker R = \ker T$.
\end{prop}

\begin{proof}[\emph{Petunjuk pembuktian}] Tinjau diagram komutatif berikut.
Gunakan \emph{teorema pemetaan terbuka}.
 \[\xy
    \scalefactor{1.4}%
    \Atrianglepair/>`>`>`<-`>/[A`B`A/\ker R`C;R`\pi`T`\wt R`\wt T]
 \endxy\]   \ns
\end{proof}

Menarik untuk diperhatikan bahwa hasil sebelumnya dapat dinyatakan dengan cara yang
sedikit memperkuat klaim keunikannya. Sebagaimana dinyatakan, kesimpulannya ialah bahwa terdapat
pemetaan linear terbatas $S$ tunggal sedemikian sehingga $SR = T$. Tanpa kerja tambahan, kita dapat merumuskan ulang
proposisi tersebut sebagai berikut.

\begin{prop}\label{prop_surj_as_factor2} Tinjau kategori ruang Banach
dan pemetaan linear kontinu. Jika $T\colon A \sto C$, $R\colon A \sto B$, $R$
surjektif, dan $\ker R \subseteq \ker T$, maka terdapat fungsi tunggal
$S$ yang memetakan $B$ ke $C$ sedemikian sehingga $S\circ R = T$. Fungsi $S$ harus
terbatas dan linear. Selain itu, fungsi tersebut injektif jika dan hanya jika $\ker R = \ker T$.
\end{prop}

Apakah versi kedua hasil ini, yang lebih terperinci, layak dinyatakan? Menurut saya, jawabannya bergantung pada
arah yang hendak dituju. Memilih versi yang lebih panjang hanya dengan alasan bahwa versi itu lebih umum
tampaknya tidak menarik. Di pihak lain, andaikan kemudian kita memerlukan hasil berikut.

\begin{prop}\label{prop_factor_blm} Jika suatu pemetaan linear terbatas $T\colon A \sto C$ antara ruang-ruang Banach dapat difaktorkan
menjadi komposisi dua fungsi $T = f \circ R$, dengan $R \in \ofml B(A,B)$ surjektif,
maka $f \colon B \sto C$ terbatas dan linear.
\end{prop}

\begin{proof}[\emph{Petunjuk pembuktian}] Untuk menerapkan~\ref{prop_surj_as_factor2}, yang diperlukan di sini hanyalah
memperhatikan bahwa $f(0) = 0$ sehingga $\ker R \subseteq \ker T$.   \ns
\end{proof}

Karena hasil sebelumnya merupakan akibat yang hampir langsung dari~\ref{prop_surj_as_factor2}, ada baiknya
versi yang lebih terperinci ini tersedia. Orang mungkin berpendapat bahwa meskipun~\ref{prop_factor_blm} barangkali tidak
sepenuhnya jelas dari pernyataan~\ref{prop_surj_as_factor}, hasil itu tentu jelas dari
pembuktiannya, jadi untuk apa memakai versi yang lebih rumit? Hal-hal semacam ini tentu merupakan soal selera.

\begin{defn}\label{defn_Bsp_exact_seq} Suatu barisan ruang Banach dan pemetaan linear kontinu
   \[\xymatrix{{\cdots}\ar[r] & B_{n-1}\ar[r]^{T_n}
          & B_n\ar[r]^{T_{n+1}} & B_{n+1}\ar[r] & {\cdots}
   }\]
dikatakan
 \index{barisan!eksak}%
 \index{eksak!barisan}%
\df{eksak di} $B_n$ jika $\ran T_n = \ker T_{n+1}$. Suatu barisan disebut \df{eksak} jika barisan tersebut
eksak di setiap ruang Banach penyusunnya. Suatu barisan ruang Banach dan
pemetaan linear kontinu berbentuk
   \begin{equation}\label{001500202i}\xymatrix{
      0 \ar[r] & A \ar[r]^S & B\ar[r]^T & C\ar[r] & 0
   }\end{equation}
 \index{barisan!eksak pendek}%
 \index{barisan eksak pendek}%
yang eksak di $A$, $B$, dan $C$ disebut \df{barisan eksak pendek}. (Di sini $\vc 0$ menyatakan ruang Banach trivial
berdimensi nol, dan panah-panah tanpa label adalah pemetaan yang jelas.)
\end{defn}

Definisi-definisi sebelumnya diberikan untuk
 \index{BAN@$\cat{BAN_\infty}$!kategori}%
 \index{kategori!$\cat{BAN_\infty}$ sebagai suatu}%
kategori $\cat{BAN_\infty}$ dari ruang Banach dan pemetaan linear kontinu. Tentu saja,
gagasan tentang suatu \emph{barisan eksak} bermakna dalam banyak situasi---misalnya, dalam
kategori aljabar Banach, aljabar-$C^*$, ruang Hilbert, ruang vektor, grup Abel,
dan modul, di antara yang lainnya.

\begin{prop}\label{00151231} Tinjau diagram berikut dalam kategori ruang Banach
dan pemetaan linear kontinu
 \[\xymatrix{
      0\ar[r] & A\ar[d]^f\ar[r]^j & B\ar[d]^g\ar[r]^k
              & C\ar@{-->}[d]^h \ar[r] & 0 \\
      0\ar[r] & A'\ar[r]_{j'} & B'\ar[r]_{k'}
      & C'\ar[r] & 0
 }\]
Jika baris-barisnya eksak dan persegi kirinya komutatif, maka terdapat pemetaan linear kontinu
$h\colon C \sto C'$ tunggal yang membuat persegi kanannya komutatif.
\end{prop}

\begin{prop}[Lema Lima]\label{00151232}
 \index{lema lima}%
Andaikan dalam diagram ruang Banach dan pemetaan linear kontinu berikut
 \[\xymatrix{
      0\ar[r] & A\ar[d]^f\ar[r]^j & B\ar[d]^g\ar[r]^k
              & C\ar[d]^h \ar[r] & 0 \\
      0\ar[r] & A'\ar[r]_{j'} & B'\ar[r]_{k'}
      & C'\ar[r] & 0
 }\]
baris-barisnya eksak dan persegi-perseginya komutatif. Maka pernyataan-pernyataan berikut berlaku.
 \begin{enumerate}[label=\rm{(\alph*)}]
  \item Jika $g$ surjektif, maka $h$ juga surjektif.
  \item Jika $f$ surjektif dan $\ran j \supseteq \ker g$, maka $h$ injektif.
  \item Jika $f$ dan $h$ surjektif, maka $g$ juga surjektif.
  \item Jika $f$ dan $h$ injektif, maka $g$ juga injektif.
 \end{enumerate}
\end{prop}

Ada banyak keuntungan bekerja dengan operator pada ruang Hilbert atau Banach yang memiliki jangkauan tertutup.
Salah satu keuntungan langsung ialah bahwa kita memperoleh kembali kesederhanaan \emph{teorema dasar aljabar linear}
(\ref{00016025}) dari versi-versi yang agak lebih rumit dalam Proposisi~\ref{prop_ker_adj_perp_ran_op}
dan Teorema~\ref{prop_FTLA_Bsp}. Di pihak lain, ada pula kerugiannya. Operator-operator itu tidak membentuk suatu kategori; komposisi
dua operator berjangkauan tertutup dapat gagal memiliki jangkauan tertutup.

\begin{prop} Jika suatu pemetaan linear terbatas $T\colon B \sto C$ antara dua ruang Banach memiliki jangkauan tertutup, maka
adjoinnya~$T^*$ juga demikian dan $\ran T^* = (\ker T)^\perp$.
\end{prop}

\begin{proof}[\emph{Petunjuk pembuktian}] Simpulkan dari Proposisi~\ref{prop_FTLA_Bsp}(d) bahwa kita hanya perlu menunjukkan bahwa
$(\ker T)^\perp \subseteq \ran T^*$. Misalkan $f$ suatu fungsional sebarang dalam $(\ker T)^\perp$. Yang harus dicari adalah suatu
fungsional $g$ dalam $C^*$ sedemikian sehingga $f = T^*g$. Untuk itu, gunakan \emph{teorema pemetaan terbuka} untuk menunjukkan bahwa pemetaan $\wt T_\ast$
dalam diagram berikut invertibel. (Pemetaan $T_\ast\colon B \sto \ran T$ dan $T$ sama persis, kecuali bahwa
kodomainnya berbeda.)

 \[\xy
    \scalefactor{1.4}%
    \Atrianglepair/>`>`>`<-`>/[B`\ran T`B/\ker T`\K;T_\ast`\pi`f`\wt{T_\ast}`\wt f]
 \endxy\]

Misalkan $g_0 := \wt f\,{\wt T_\ast}^{-1}$. Gunakan \emph{teorema Hahn--Banach} untuk menghasilkan~$g$ yang diinginkan.   \ns
\end{proof}

\begin{prop} Dalam kategori $\cat{BAN_\infty}$, jika barisan $\xymatrix{A \ar[r]^S & B \ar[r]^T & C}$ eksak di $B$
dan $T$ memiliki jangkauan tertutup, maka $\xymatrix{C^* \ar[r]^{T^*} & B^* \ar[r]^{S^*} & A^*}$ eksak di~$B^*$.
\end{prop}

\begin{defn} Suatu funktor dari kategori ruang Banach dan pemetaan linear kontinu ke dirinya sendiri disebut
 \index{eksak!funktor}%
 \index{funktor!eksak}%
\df{eksak} jika funktor tersebut membawa barisan eksak pendek ke barisan eksak pendek.
\end{defn}

\begin{exam}\label{exam_dualfunctor_exact} Funktor $B \mapsto B^*$, $T \mapsto T^*$ (lihat~\ref{exam_dual_functor})
dari kategori ruang Banach dan pemetaan linear kontinu ke dirinya sendiri bersifat eksak.
\end{exam}

\begin{exam} Misalkan $B$ suatu ruang Banach dan $\sbsb{\tau}B\colon x \mapsto x^{**}$ pembenaman alami dari $B$ ke~$B^{**}$
(lihat Definisi~\ref{defn_nat_embed}). Maka ruang hasil bagi $B^{**}/\ran\sbsb{\tau}B$ merupakan ruang Banach.
\end{exam}

\begin{notn} Kita akan menyatakan dengan $\clo B$ ruang hasil bagi
$B^{**}/\ran \sbsb{\tau}B$ dalam contoh sebelumnya.
\end{notn}

\begin{prop} Jika $T\colon B \sto C$ suatu pemetaan linear kontinu antara ruang-ruang Banach, maka terdapat pemetaan linear kontinu tunggal
$\clo T\colon \clo B \sto \clo C$ sedemikian sehingga $\clo T(\phi + \ran\sbsb{\tau}B) = (T^{**}\phi) + \ran\sbsb{\tau}C$ untuk setiap $\phi \in B^{**}$.
\end{prop}

\begin{proof}[\emph{Petunjuk pembuktian}] Terapkan~\ref{00151231} pada diagram berikut:
 \begin{equation}\label{eqn_exactCD_Bbar}\xymatrix{
      0\ar[r] & B\ar[d]^T\ar[r]^{\sbsb{\tau}B} & B^{**}\ar[d]^{T^{**}}\ar[r]^{\sbsb{\pi}B}
              & \clo B\ar@{-->}[d]^{\clo T} \ar[r] & 0 \\
      0\ar[r] & C\ar[r]_{\sbsb{\tau}C} & C^{**}\ar[r]_{\sbsb{\pi}C} & \clo C\ar[r] & 0
 }\end{equation}  \ns
\end{proof}

\begin{exam}\label{exam_bar_functor} Pasangan pemetaan $B \mapsto \clo B$, $T \mapsto \clo T$ dari
kategori $\cat{BAN_\infty}$ ke dirinya sendiri merupakan suatu
 \index{<unopu@$B \mapsto \clo B$, $T \mapsto \clo T$ (funktor)}%
 \index{funktor!dari suatu ruang Banach $B$ ke $\clo B$}%
funktor kovarian eksak.
\end{exam}

\begin{prop} Dalam kategori ruang Banach dan pemetaan linear kontinu, jika
barisan
 \[\begin{CD}
  0 \longrightarrow A @> S >> B @> T >> C \longrightarrow 0
  \end{CD}\]
eksak, maka diagram berikut komutatif dan memiliki baris-baris
serta kolom-kolom yang eksak.
 \[\bfig
    \iiixiii{'7777}[A`A^{**}`\overline A`B`B^{**}`\overline B`C`C^{**}`\overline C;%
``````S`{S^{**}}`{\overline S}`T`{T^{**}}`{\overline T}]
  \efig\]
\end{prop}

\begin{cor} Setiap subruang dari suatu ruang Banach refleksif bersifat refleksif.
\end{cor}

\begin{cor} Jika $M$ suatu subruang dari ruang Banach refleksif~$B$, maka $B/M$ refleksif.
\end{cor}

\begin{cor} Misalkan $M$ suatu subruang dari ruang Banach~$B$. Jika $M$ dan $B/M$ keduanya refleksif,
maka $B$ juga refleksif.
\end{cor}



















\begin{defn} Misalkan $T \in \ofml L(B,C)$, dengan $B$ dan $C$ ruang Banach. Definisikan
$\coker T$, yaitu
 \index{kokernel}%
 \index{kokernel@$\coker T$ (kokernel dari $T$)}%
\df{kokernel} dari~$T$, sebagai $C/\ran T$.
\end{defn}



Menentukan secara eksplisit dual dari suatu ruang Banach sering kali agak menantang. Berikut diberikan beberapa
contoh penting. Namun, pertama-tama perlu disampaikan satu peringatan. Notasi untuk objek-objek dalam ruang linear bernorma
yang terdiri atas barisan dapat menyesatkan. Ketika membahas kelengkapan ruang-ruang semacam itu, kita berhadapan dengan barisan dari
barisan bilangan kompleks. Demi memudahkan pembaca (dan mungkin juga diri Anda sendiri sebulan kemudian,
ketika Anda mencoba membaca sesuatu yang Anda tulis!), sebaiknya bedakan secara notasi antara barisan bilangan kompleks
dan barisan dari barisan. Ada banyak cara untuk melakukannya: subskrip ganda, subskrip dan
superskrip, notasi fungsional untuk barisan kompleks, \emph{dan sebagainya}.

Notasi bermasalah lain yang muncul dalam ruang barisan adalah $\sum_{k=1}^\infty a_k \vc e^k$, dengan $\vc e^k$
merupakan barisan yang bernilai $1$ pada koordinat ke-$k$ dan $0$ di semua koordinat lainnya. (Ingat bahwa dalam
Contoh~\ref{C063149}, kita menggunakan vektor-vektor ini sebagai basis ortonormal biasa (atau standar) untuk ruang Hilbert
$l_2$.) Tampaknya hampir tidak ada kebingungan yang timbul jika $\sum_{k=1}^\infty a_k \vc e^k$ dipandang
hanya sebagai singkatan untuk barisan bilangan kompleks $(a_k) = (a_1,a_2,a_3,\dots)$. Namun, jika seseorang hendak
memperlakukannya secara harfiah sebagai deret tak hingga, diperlukan kehati-hatian. Dalam ruang $\sbsb c0$ dan $l_1$, tidak ada
masalah. Deret itu memang merupakan limit, ketika $n$ membesar, dari $\sum_{k=1}^n a_k \vc e^k$. (Mengapa?) Namun, dalam ruang
$c$ dan $l_\infty$, misalnya, terjadi kegagalan yang sangat buruk. Untuk melihat apa yang dapat terjadi, tinjau barisan konstan
$(1,1,1,\dots)$ (yang merupakan anggota $c$ maupun $l_\infty$). Maka
$\lim_{n \sto\infty}\norm{(1,1,1,\dots) - \sum_{k=1}^n \vc e^k}$ gagal mendekati nol dalam kedua kasus tersebut.
Dalam $c$ maupun $l_\infty$, limitnya adalah~$1$.

\begin{exam}\label{exam_dual_C0} Ruang linear bernorma $\sbsb c0$ (lihat~\ref{C0231303})
 \index{c@$c_0$!sebagai ruang Banach}%
 \index{Banach!ruang!$c_0$ sebagai suatu}%
 \index{l@$l_1$!sebagai dual dari $c_0$}%
 \index{dual!dari $c_0$}%
merupakan ruang Banach dan dualnya adalah $l_1$ (lihat~\ref{C023191}).
\end{exam}

\begin{proof}[\emph{Petunjuk pembuktian}] Tinjau pemetaan
$T\colon {\sbsb c0}^* \sto l_1\colon f \mapsto \sum_{k=1}^\infty f(\vc e^k)\vc e^k$.
(Jangan lupa menunjukkan bahwa $T$ benar-benar memetakan ke~$l_1$.) Periksa bahwa $T$ merupakan isomorfisme isometrik.
Untuk membuktikan bahwa $T$ surjektif, mulailah dengan suatu vektor $b \in l_1$ dan tinjau fungsi
$f\colon \sbsb c0 \sto \C\colon a \mapsto \sum_{k=1}^\infty a_kb_k$.    \ns
\end{proof}

\begin{cor} Ruang $l_1$ merupakan suatu
 \index{l@$l_1$!sebagai ruang Banach}%
 \index{Banach!ruang!$l_1$ sebagai suatu}%
ruang Banach.
\end{cor}

Menarik bahwa ternyata ruang $\sbsb c0$ dan $c$ tidak isomorfik secara isometrik,
meskipun ruang-ruang dualnya isomorfik secara isometrik. Perhitungan dual dari $c$ secara teknis sedikit lebih rumit daripada
perhitungan~${\sbsb c0}^{\!*}$. Mungkin berguna untuk meninjau \emph{basis Schauder} bagi kedua ruang ini.

\begin{defn} Suatu barisan $(\vc e^k)$ dari vektor-vektor dalam suatu ruang Banach $B$ disebut
 \index{basis Schauder}%
 \index{basis!Schauder}%
\df{basis Schauder} bagi ruang tersebut jika untuk setiap $x \in B$ terdapat barisan skalar $(\alpha_k)$ yang unik
sedemikian sehingga $x = \sum_{k=1}^\infty \alpha_k \vc e^k$.
\end{defn}

Dua catatan mengenai definisi sebelumnya: Kata ``unik'' sangat penting; dan tidak benar bahwa
setiap ruang Banach memiliki basis Schauder.

\begin{exam} Suatu basis ortonormal dalam ruang Hilbert separabel merupakan basis Schauder bagi ruang tersebut. (Lihat
Proposisi~\ref{prop_unique_onbasis} dan~\ref{prop_Schauder_basis_Hsp}.)
\end{exam}

\begin{exam}\label{exam_dual_c0} Kita memilih
 \index{biasa!basis Schauder untuk~$\sbsb c0$}%
 \index{Schauder!basis!biasa (atau standar) untuk~$\sbsb c0$}%
 \index{basis!standar (atau biasa)}%
vektor-vektor yang disebut \df{standar} (atau
 \index{standar!basis Schauder untuk $\sbsb c0$}%
\df{biasa}) sebagai \df{vektor-vektor basis} untuk $\sbsb c0$ dengan cara yang sama seperti untuk~$l_2$ (lihat~\ref{C063149}).
Untuk setiap $n \in \N$, misalkan $\vc e^n$ adalah barisan dalam $l_2$ yang koordinat ke-$n$-nya bernilai $1$
dan semua koordinat lainnya bernilai~$0$. Maka $\{\vc e^n\colon n \in \N\}$ merupakan basis Schauder
bagi~$\sbsb c0$.
\end{exam}

Sepintas lalu, ruang $c$ dari semua barisan konvergen mungkin tampak jauh lebih besar daripada $\sbsb c0$,
yang hanya memuat barisan-barisan yang konvergen ke nol. Lagi pula, untuk setiap vektor
$x$ dalam $\sbsb c0$ terdapat tak terhitung banyaknya vektor berbeda dalam $c$ (cukup tambahkan sebarang bilangan kompleks
tak nol pada setiap suku barisan~$x$). Mungkinkah $c$, yang memuat tak terhitung banyaknya
salinan dari koleksi anggota~$\sbsb c0$ yang tak terhitung, memiliki basis Schauder (yang mesti terhitung)?
Setelah direnungkan, jawabannya adalah \emph{ya}. Kita hanya perlu menambahkan satu elemen
pada basis bagi~$\sbsb c0$ untuk menangani translasi.

\begin{exam} Misalkan $\vc e^1$, $\vc e^2$, \dots seperti dalam Contoh~\ref{exam_dual_c0} sebelumnya. Selain itu,
nyatakan dengan $\vc 1$ barisan konstan $(1,1,1,\dots)$ dalam~$c$. Maka $(\vc 1, \vc e^1, \vc e^2, \dots)$
merupakan basis Schauder bagi~$c$.
\end{exam}

\begin{exam}\label{exam_dual_c} Ruang linear bernorma $c$ (lihat~\ref{C0231302})
 \index{c@$c$!sebagai ruang Banach}%
 \index{Banach!ruang!$c$ sebagai suatu}%
 \index{l@$l_1$!sebagai dual dari $c$}%
 \index{dual!dari $c$}%
merupakan ruang Banach dan dualnya adalah $l_1$ (lihat~\ref{C023191}).
\end{exam}

\begin{proof}[\emph{Petunjuk pembuktian}]
\leavevmode\par\noindent
Pertama-tama, perhatikan bahwa untuk setiap $f \in c^*$, barisan
$\bigl(f(\vc e^k)\bigr)$ dapat dijumlahkan secara mutlak dan, sebagai akibatnya, $\sum_{k=1}^\infty a_k f(\vc e^k)$
konvergen untuk setiap barisan $a = (a_1,a_2,\dots)$ dalam~$c$. Perhatikan pula bahwa setiap $a \in c$ dapat dituliskan sebagai
$a = a_\infty \vc 1 + \sum_{k=1}^\infty (a_k -a_\infty)\vc e^k$, dengan $a_\infty := \lim_{k \sto \infty}a_k$
dan $\vc 1 := (1,1,1,\dots)$. Gunakan hal ini untuk membuktikan bahwa untuk setiap $f \in c^*$,
   \[ \abs{f(a)} \le \biggl[ \abs{\beta_f} + \sum_{k=1}^\infty \abs{f(\vc e^k)} \biggr] {\norm a}_\infty \]
dengan $\beta_f := f(\vc 1) - \sum_{k=1}^\infty f(\vc e^k)$. Untuk $t = (t_0,t_1,t_2,\dots) \in l_1$ dan
$a = (a_1,a_2,a_3,\dots) \in c$, definisikan
   \[ S_t(a) := t_0 a_\infty + \sum_{k=1}^\infty t_ka_k \]
dan periksa bahwa $S\colon t \mapsto S_t$ merupakan pemetaan linear kontinu dari $l_1$ ke~$c^*$.
Selanjutnya, untuk $f \in c^*$, definisikan
   \[ Tf := (\beta_f, f(\vc e^1), f(\vc e^2), \dots) \]
dan tunjukkan bahwa $T\colon f \mapsto Tf$ merupakan isometri linear dari $c^*$ ke~$l_1$. Sekarang, untuk melihat bahwa $c^*$ dan
$l_1$ isomorfik secara isometrik, cukup ditunjukkan bahwa $ST = \id{c^*}$ dan $TS = \id{l_1}$.  \ns
\end{proof}












\begin{defn} Suatu ruang topologis disebut
 \index{lokal!kompak}%
 \index{kompak!lokal}%
\df{kompak lokal} jika setiap titik dalam ruang tersebut memiliki lingkungan kompak.
\end{defn}

\begin{notn} Misalkan $X$ suatu ruang kompak lokal. Suatu fungsi bernilai skalar $f$ pada $X$ termasuk dalam $\fml C_0(X)$ jika
 \index{C@$\fml C_0(X)$!ruang fungsi kontinu yang lenyap di tak hingga}%
fungsi tersebut kontinu dan lenyap di tak hingga (yakni, jika untuk setiap bilangan $\epsilon > 0$ terdapat subhimpunan kompak
dari $X$ yang di luarnya berlaku $\abs{f(x)} < \epsilon$).
\end{notn}

\begin{exam}\label{exam_C0X} Jika $X$ suatu ruang Hausdorff kompak lokal, maka, terhadap operasi aljabar titik demi titik dan norma seragam,
 \index{Banach!ruang!$\fml C_0(X)$ sebagai suatu}%
 \index{C@$\fml C_0(X)$!sebagai ruang Banach}%
$\fml C_0(X)$ merupakan ruang Banach.
\end{exam}

Kita mengingat kembali salah satu teorema terpenting dari analisis real.

\begin{thm}[Representasi Riesz]\label{C057951} Misalkan $X$ suatu ruang Hausdorff kompak lokal.
 \index{Riesz!teorema representasi}%
Jika $\phi$ suatu fungsional linear terbatas pada $\fml C_0(X)$, maka terdapat ukuran Borel kompleks reguler
$\mu$ yang unik pada $X$ sedemikian sehingga
  \[ \phi(f) = \int_X f\,d\mu \]
untuk setiap $f \in \fml C_0(X)$ dan $\norm\mu = \norm\phi$.
\end{thm}

\begin{proof} Teorema ini dan akibat berikutnya disajikan dengan baik dalam Lampiran C pada buku analisis
fungsional karya Conway~\cite{Conway:1990}. Lihat juga~\cite{Folland:1984}, Teorema 7.17; \cite{HewittS:1965},
Teorema 20.47 dan 20.48; \cite{McDonaldW:1999}, Teorema 9.16; dan \cite{Rudin:1987}, Teorema 6.19.  \ns
\end{proof}

\begin{cor}\label{C057952} Jika $X$ suatu ruang Hausdorff kompak lokal, maka $\bigl(\fml C_0(X)\,\bigr)^*$
 \index{dual!dari $\fml C_0(X)$}%
 \index{C@$\fml C_0(X)$!dual dari}%
 \index{M@$\textbf{M}(X)$!sebagai dual dari $\fml C_0(X)$}%
isomorfik secara isometrik dengan~$\textbf{M}(X)$,
 \index{M@$\textbf{M}(X)$!ukuran Borel kompleks reguler pada~$X$}%
ruang Banach dari semua ukuran Borel kompleks reguler pada~$X$.
\end{cor}

\begin{cor} Jika $X$ suatu ruang Hausdorff kompak lokal, maka $\textbf{M}(X)$
 \index{M@$\textbf{M}(X)$!sebagai ruang Banach}%
 \index{Banach!ruang!$\textbf{M}(X)$ sebagai suatu}%
merupakan ruang Banach.
\end{cor}










Dalam~\ref{defn_nls_adjoint}, kita mendefinisikan adjoin dari suatu operator linear antara ruang-ruang Banach,
dan dalam~\ref{def000926}, kita mendefinisikan adjoin dari suatu operator ruang Hilbert.
Dalam kasus operator pada ruang Hilbert (yang juga merupakan ruang Banach),
wajar untuk menanyakan hubungan apa, jika ada, di antara kedua ``adjoin'' ini.
Keduanya tentu tidak dapat sama, karena yang pertama bekerja di antara ruang-ruang dual, sedangkan yang kedua
di antara ruang-ruang semula. Dalam proposisi berikut, kita menggunakan antiisomorfisme
$\psi$ yang didefinisikan dalam~\ref{000341} (lihat~\ref{0022057}) untuk menunjukkan bahwa kedua adjoin itu
``pada hakikatnya'' sama.

\begin{prop} Misalkan $T$ suatu operator pada ruang Hilbert $H$ dan $\psi$ antiisomorfisme yang
didefinisikan dalam~\ref{000341}. Jika kita (untuk sementara) menyatakan dual ruang Banach
dari $T$ dengan $T'\colon H^* \sto H^*$, maka $T' = \psi\, T^* \psi^{-1}$. Artinya,
diagram berikut komutatif.
   \[ \xy
          \Square/{->}`{<-}`{<-}`{->}/[H^*`H^*`H`H;T'`\psi`\psi`T^*]
      \endxy \]
\end{prop}

\begin{defn}\label{0022091} Misalkan $A$ suatu subhimpunan dari ruang Banach $B$. Maka
 \index{anihilator}%
\df{anihilator} dari $A$, yang dinotasikan
 \index{<perp@$A^\perp$ (anihilator dari suatu himpunan)}%
dengan~$A^\perp$, adalah $\{f \in B^*\colon f(a) = 0 \text{ untuk setiap $a \in A$}\}$.
\end{defn}

Konflik notasi antara \ref{0001511} dan \ref{0022091} dapat
menimbulkan kebingungan. Untuk melihat bahwa dalam ruang Hilbert, anihilator suatu subruang dan
komplemen ortogonal subruang itu ``pada hakikatnya'' merupakan hal yang sama, gunakan
antiisomorfisme isometrik $\psi$ antara $H$ dan $H^*$ yang dibahas dalam~\ref{0022057} untuk mengidentifikasi keduanya.

\begin{prop} Misalkan $M$ suatu subruang dari ruang Hilbert $H$. Jika kita (untuk sementara) menyatakan
anihilator dari $M$ dengan~$M^a$, maka $M^a = \psi^{\sto}\bigl(M^\perp\bigr)$.
\end{prop}
\section{Proyeksi dan Subruang Terkomplemen}
\begin{defn}
 Pemetaan linear terbatas dari suatu ruang Banach ke dirinya sendiri disebut
 \index{proyeksi!di ruang Banach}%
 \index{operator!proyeksi}%
\df{proyeksi} jika bersifat idempoten. Jika $E$ adalah pemetaan seperti itu, kita menyebutnya
proyeksi \df{pada $\ran E$ sepanjang $\ker E$}.
\end{defn}

Perhatikan bahwa ini sudah pasti \emph{bukan} hal yang sama dengan proyeksi di
ruang Hilbert. Proyeksi ruang Hilbert adalah \emph{proyeksi ortogonal}; yakni,
proyeksi itu swaadjoin sekaligus idempoten. Jangkauan suatu proyeksi ruang Hilbert
tegak lurus terhadap kernelnya. Karena istilah ``ortogonal'', ``swaadjoin'', dan
``tegak lurus'' tidak bermakna di ruang Banach, kita tidak akan mengharapkan kedua pengertian
``proyeksi'' itu identik. Namun, dalam kedua kasus tersebut proyeksi \emph{memang}
terbatas, linear, dan idempoten.

\begin{defn} Misalkan $M$ subruang linear tertutup dari suatu ruang Banach~$B$. Jika terdapat
subruang linear tertutup $N$ dari $B$ sedemikian sehingga $B = M \oplus N$, maka kita mengatakan bahwa $M$ adalah
 \index{subruang terkomplemen}%
 \index{subruang!terkomplemen}%
subruang \df{terkomplemen}, bahwa $N$ adalah
 \index{komplemen!dari subruang Banach}%
\df{komplemen (ruang Banach)}-nya, dan bahwa subruang $M$ dan $N$ bersifat
\df{komplementer}.
\end{defn}

\begin{prop}\label{C069814} Misalkan $M$ dan $N$ subruang komplementer dari suatu ruang Banach~$B$. Jelas bahwa setiap
elemen di $B$ dapat ditulis secara tunggal dalam bentuk $m + n$ untuk suatu $m \in M$ dan $n \in
N$. Definisikan pemetaan $E \colon B \sto B$ dengan $E(m + n) = m$ untuk $m \in M$ dan $n \in N$.
Maka $E$ adalah proyeksi pada $M$ sepanjang~$N$.
\end{prop}

\begin{proof}[\emph{Petunjuk pembuktian}] Gunakan~\ref{prop_cont_idem_op}.  \ns
\end{proof}

\begin{prop}\label{C069817} Jika $E$ adalah proyeksi pada suatu ruang Banach $B$, maka kernel dan jangkauannya merupakan
subruang-subruang komplementer dari~$B$.
\end{prop}

Berikut ini merupakan akibat langsung dari dua proposisi sebelumnya.

\begin{cor}\label{cor_compl_ran_proj} Suatu subruang dari ruang Banach terkomplemen jika dan hanya jika subruang itu
merupakan jangkauan suatu proyeksi.
\end{cor}

\begin{prop} Misalkan $B$ ruang Banach. Jika $E \colon B \sto B$ adalah proyeksi ke suatu subruang
$M$ sepanjang subruang $N$, maka $I - E$ adalah proyeksi ke $N$ sepanjang~$M$.
\end{prop}

\begin{prop}\label{C069824} Jika $M$ dan $N$ merupakan subruang komplementer dari suatu ruang Banach $B$, maka $B/M$ dan
$N$ isomorfik.
\end{prop}

\begin{prop}\label{prop_linv_Bsp} Misalkan $T \colon B \sto C$ pemetaan linear terbatas di antara ruang-ruang Banach.
 \index{kiri!invers!operator ruang Banach}%
Jika $T$ injektif dan $\ran T$ terkomplemen, maka $T$ mempunyai invers kiri.
\end{prop}

\begin{proof}[\emph{Petunjuk pembuktian}] Untuk petunjuk ini kita memperkenalkan notasi (yang agak rumit).
Jika $f\colon X \sto Y$ adalah fungsi di antara dua himpunan, kita definisikan
$f_\bullet \colon X \sto \ran f\colon x \mapsto f(x)$. Artinya, $f_\bullet$ tepat merupakan fungsi yang sama
dengan $f$, kecuali bahwa kodomainnya adalah jangkauan~$f$.

Karena $\ran T$ terkomplemen, terdapat proyeksi $E$ dari $C$ ke~$\ran T$
(menurut~\ref{cor_compl_ran_proj}). Tunjukkan bahwa $T_\bullet$ mempunyai invers kontinu.
Misalkan $S = {T_\bullet}^{-1}E_\bullet$.    \ns
\end{proof}

\begin{prop}\label{prop_rinv_Bsp} Misalkan $S \colon C \sto B$ pemetaan linear terbatas di antara ruang-ruang Banach.
 \index{kanan!invers!operator ruang Banach}%
Jika $S$ surjektif dan $\ker S$ terkomplemen, maka $S$ mempunyai invers kanan.
\end{prop}

\begin{proof}[\emph{Petunjuk pembuktian}] Karena $\ker S$ terkomplemen, terdapat $N \preccurlyeq C$ sedemikian sehingga
$C = N \oplus \ker S$. Tunjukkan bahwa pembatasan $S$ ke $N$ mempunyai invers kontinu dan komposisikan
invers ini dengan pemetaan inklusi dari $N$ ke~$C$.   \ns
\end{proof}

Hasil berikut merupakan konvers dari kedua proposisi~\ref{prop_linv_Bsp} dan~\ref{prop_rinv_Bsp}.

\begin{prop}\label{C069831} Misalkan $T \colon B \sto C$ dan $S\colon C \sto B$ pemetaan linear terbatas di antara ruang-ruang
Banach. Jika $ST = \id B$, maka
 \begin{enumerate}[label=\rm{(\alph*)}]
  \item[(a)] $S$ surjektif;
  \item[(b)] $T$ injektif;
  \item[(c)] $TS$ merupakan proyeksi pada $\ran T$ sepanjang~$\ker S$; dan
  \item[(d)] $C = \ran T \oplus \ker S$.
 \end{enumerate}
\end{prop}

\begin{prop} Jika dalam kategori $\cat{BAN_\infty}$ barisan
 \[ \begin{CD} 0 @>>> A @>j>> B @>f>> C @>>> 0
    \end{CD}\]
 dan
 \[ \begin{CD} 0 @>>> C @>g>> B @>p>> D @>>> 0
    \end{CD}\]
eksak dan jika $fg = \id{C}$, maka $pj$ adalah isomorfisme dari $A$ pada~$D$.
\end{prop}

\begin{proof}[\emph{Petunjuk pembuktian}] Gunakan \ref{C069831} untuk melihat bahwa $B = \ker p \oplus \ran j$. Tinjau pemetaan
$\bigl.p\bigr|_{\ran j}\,j_\bullet$ (notasi untuk $j_\bullet$ seperti dalam petunjuk untuk~\ref{prop_linv_Bsp}). \ns
\end{proof}

\begin{prop}\label{prop_merge_exactCD} Jika dalam kategori $\cat{BAN_\infty}$ diagram
 \begin{equation} \begin{CD} 0 @>>> A @>j>> B @>f>> C @>>> 0  \\
      @.   @VSVV   @VVUV   @.    @.    \\
 0 @>>> A' @>>j'> B' @>>f'> C' @>>> 0
    \end{CD}\end{equation}
dan
 \[ \begin{CD} 0 @>>> C @>g>> B @>p>> D @>>> 0   \\
      @.     @.   @VUVV   @VVTV     @.  \\
  0 @>>> C' @>>g'> B' @>>p'> D' @>>> 0
    \end{CD}\]
komutatif, jika baris-barisnya eksak, jika $fg = \id{C}$ dan jika $f'g' = \id{C'}$, maka diagram
 \[ \begin{CD}  A  @>pj>>  D   \\
       @VSVV   @VVTV  \\
     A'  @>>p'j'> D'
    \end{CD}\]
komutatif dan pemetaan $pj$ serta $p'j'$ merupakan isomorfisme.
\end{prop}

Dalam~\ref{exam_dual_functor} dan~\ref{exam_dualfunctor_exact} kita melihat bahwa dalam kategori ruang Banach
dan pemetaan linear kontinu, funktor dualitas $B \mapsto B^*$, $T \mapsto T^*$ merupakan funktor kontravarian eksak.
Dan dalam~\ref{exam_bar_functor} kita melihat bahwa dalam kategori yang sama pasangan pemetaan $B \mapsto \clo B$,
$T \mapsto \clo T$ merupakan funktor kovarian eksak. Dengan mengomposisikan kedua funktor ini diperoleh dua funktor
kontravarian eksak baru, yang akan kita lambangkan dengan $F$ dan $G$:
  \[ \xymatrix{B\ar[r]^F &{\clo B}^{\,*}}\,,\,\xymatrix{T\ar[r]^F &{\clo T}^{\,*}} \qquad \text{dan} \qquad
           \xymatrix{B\ar[r]^G &\clo{B^*}}\,,\,\xymatrix{T\ar[r]^G &\clo{T^*}}  \]
Wajar untuk menanyakan hubungan antara ${\clo B}^{\,*}$ dan $\clo{B^*}$. Proposisi berikut
memberi tahu kita bukan hanya bahwa kedua ruang ini selalu isomorfik, melainkan bahwa keduanya isomorfik secara alami. Artinya,
 \index{alami!ekuivalensi!antara ${\clo B}^{\,*}$ dan $\clo{B^*}$}%
funktor $F$ dan $G$ ekuivalen secara alami.

\begin{prop} Jika $T \in \ofml B(B,C)$, dengan $B$ dan $C$ ruang Banach, maka terdapat isomorfisme
$\sbsb{\phi}B \colon {\clo B}^{\,*} \sto \clo{B^*}$ dan $\sbsb{\phi}C \colon {\clo C}^{\,*} \sto \clo{C^*}$ sedemikian sehingga
   \[ {\clo T}^{\,*} = {\sbsb{\phi}B}^{\!\!\!-1}\,\, \clo{T^*}\,\,\sbsb{\phi}C\,. \]
\end{prop}

\begin{proof}[\emph{Petunjuk pembuktian}] Karena funktor dualitas eksak (lihat~\ref{exam_dualfunctor_exact}), kita dapat
mengambil dual dari diagram komutatif~\eqref{eqn_exactCD_Bbar} untuk memperoleh diagram komutatif kedua dengan baris-baris
eksak. Sekarang tuliskan diagram komutatif ketiga dengan mengganti setiap kemunculan $B$ dalam diagram~\eqref{eqn_exactCD_Bbar}
dengan~$B^*$. Periksa bahwa ${\sbsb{\tau}B}^{\!\!*}\,\sbsb{\tau}{B^*} = \id{B^*}$ sehingga Anda dapat menggunakan proposisi~\ref{prop_merge_exactCD}
pada dua diagram terakhir dan dengan demikian memperoleh diagram komutatif
     \[ \xy
          \Square/{<-}`{->}`{->}`{<-}/[{\clo B}^{\,*}`{\clo C}^{\,*}`\clo{B^*}`\clo{C^*};{\clo T}^{\,*}`\sbsb{\phi}B`\sbsb{\phi}C`\clo{T^*}]
      \endxy \]
dengan $\sbsb{\phi}B := \sbsb{\pi}{B^*}\,{\sbsb{\pi}B}^{\!\!*}$ suatu isomorfisme.   \ns
\end{proof}

\begin{cor} Suatu ruang Banach refleksif jika dan hanya jika dualnya refleksif.
\end{cor}

\begin{prop}\label{prop_fd_subsp_complm} Setiap subruang berdimensi hingga dari suatu ruang Banach terkomplemen.
\end{prop}

\begin{proof}[\emph{Petunjuk pembuktian}] Misalkan $F$ subruang berdimensi hingga dari suatu ruang Banach $B$ dan misalkan
$\{\vc e^1, \dots, \vc e^n\}$ basis Hamel untuk~$F$. Maka untuk setiap $x \in F$ terdapat skalar-skalar tunggal
$\alpha_1(x)$, \dots , $\alpha_n(x)$ sedemikian sehingga $x = \sum_{k=1}^n \alpha_k(x)\vc e^k$. Untuk setiap $1 \le k \le n$
verifikasi bahwa $\alpha_k$ merupakan fungsional linear terbatas, yang dapat diperluas ke seluruh~$B$. Misalkan $N$ adalah
irisan kernel-kernel dari perluasan ini dan tunjukkan bahwa $B = F \oplus N$.  \ns
\end{proof}

\begin{defn} Ingat bahwa dalam~\ref{000082} kita mendefinisikan \emph{kodimensi} suatu subruang dari ruang vektor
sebagai dimensi komplemen ruang vektornya (yang selalu ada menurut proposisi~\ref{0000824}). Kita menggunakan
istilah yang sama untuk ruang Banach:
 \index{kodimensi}%
\df{kodimensi} dari sebarang \emph{subruang vektor} suatu ruang Banach adalah dimensi \emph{komplemen ruang
vektornya}.
\end{defn}

Untuk korolar berikut (dari proposisi~\ref{prop_fd_subsp_complm}), ingatlah
konvensi~\ref{conv_subsp_Bsp} tentang penggunaan kata ``subruang'' dalam konteks ruang Banach.

\begin{cor} Dalam suatu ruang Banach, setiap subruang berkodimensi hingga terkomplemen.
\end{cor}





















Misalkan $M$ subruang tak nol dari suatu ruang Banach~$B$. Tinjau dua sifat berikut dari pasangan $B$ dan~~$M$.

\textbf{Sifat P:} \emph{Terdapat proyeksi bernorma $1$ dari $B$ ke $M$.}

\textbf{Sifat E:} \emph{Untuk setiap ruang Banach $C$, setiap pemetaan linear terbatas dari $M$ ke $C$ dapat diperluas menjadi pemetaan
linear terbatas dari $B$ ke $C$ tanpa memperbesar normanya.}

\begin{prop} Untuk subruang tak nol $M$ dari suatu ruang Banach $B$, sifat \textbf{P} dan \textbf{E} ekuivalen.
\end{prop}

\begin{proof}[\emph{Petunjuk pembuktian}] Untuk menunjukkan bahwa \textbf{E} mengakibatkan \textbf{P}, mulailah dengan memperluas pemetaan identitas
pada $M$ ke seluruh~$B$.   \ns
\end{proof}

















\section{Prinsip Keterbatasan Seragam}
\begin{defn} Misalkan $S$ suatu himpunan dan $V$ ruang linear bernorma. Suatu keluarga $\fml F$ dari
fungsi-fungsi dari $S$ ke $V$ disebut
 \index{titik demi titik!terbatas}%
 \index{terbatas!titik demi titik}%
\df{terbatas titik demi titik} jika untuk setiap $x \in S$ terdapat konstanta $M_x > 0$ sedemikian sehingga
$\norm{f(x)} \le M_x$ untuk setiap $f \in \fml F$.
\end{defn}

\begin{defn} Misalkan $V$ dan $W$ ruang linear bernorma. Suatu keluarga $\fml T$ dari pemetaan linear terbatas
dari $V$ ke $W$ disebut
 \index{seragam!terbatas}%
 \index{terbatas!seragam}%
\df{terbatas seragam} jika terdapat $M > 0$ sedemikian sehingga $\norm T \le M$ untuk setiap $T \in
\fml T$.
\end{defn}

\begin{exer}\label{C070114} Misalkan $V$ dan $W$ ruang linear bernorma. Jika suatu keluarga $\fml T$ dari
pemetaan linear terbatas dari $V$ ke $W$ terbatas seragam, maka keluarga itu terbatas titik demi titik.
\end{exer}

\emph{Prinsip keterbatasan seragam} adalah pernyataan bahwa konvers
dari~\ref{C070114} berlaku apabila $V$ lengkap. Untuk membuktikannya kita akan menggunakan
analog ``lokal'' dari pernyataan ini untuk fungsi-fungsi kontinu pada ruang metrik lengkap. Secara kasar, analog ini mengatakan bahwa jika
suatu keluarga fungsi bernilai real pada ruang metrik lengkap terbatas titik demi titik, maka keluarga itu terbatas
seragam pada suatu subhimpunan terbuka tak kosong dari ruang tersebut.

\begin{prop}\label{C070117} Misalkan $M$ ruang metrik lengkap yang tak kosong dan $\fml F \subseteq \fml C(M,\R)$.
Jika untuk setiap $x \in M$ terdapat konstanta $N_x > 0$ sedemikian sehingga $\abs{f(x)} \le N_x$ untuk
setiap $f \in \fml F$, maka terdapat himpunan terbuka tak kosong $U$ di $M$ dan bilangan $N > 0$ sedemikian
sehingga $\abs{f(u)} \le N$ untuk setiap $f \in \fml F$ dan setiap $u \in U$.
\end{prop}

\begin{proof}[\emph{Petunjuk pembuktian}] Untuk setiap bilangan asli $k$, misalkan $A_k =
\bigcap\{f^{\gets}\bigl(\,[-k,k]\,\bigr)\colon f \in \fml F\}$. Gunakan versi yang sama dari
\emph{teorema kategori Baire} yang kita gunakan dalam membuktikan \emph{teorema pemetaan terbuka} (\ref{C069414})
untuk menyimpulkan bahwa untuk sedikitnya satu $n \in \N$, himpunan $A_n$ mempunyai interior
tak kosong. Misalkan $U = \intr{A_n}$.   \ns
\end{proof}

\begin{thm}[Prinsip keterbatasan seragam]\label{thm_PUB} Misalkan $B$ ruang Banach dan $W$ ruang linear bernorma.
 \index{prinsip!keterbatasan seragam}%
 \index{keterbatasan!prinsip seragam}%
 \index{seragam!keterbatasan!prinsip}%
Jika suatu keluarga $\ofml T$ dari pemetaan linear terbatas dari $B$ ke $W$ terbatas titik demi
titik, maka keluarga itu terbatas seragam.
\end{thm}

\begin{proof}[\emph{Petunjuk pembuktian}] Untuk setiap $T \in \ofml T$, definisikan
  \[ \sbsb fT\colon B \sto \R\colon x \mapsto \norm{Tx}. \]
Gunakan proposisi sebelumnya untuk menunjukkan bahwa terdapat subhimpunan terbuka tak kosong $U$ dari $B$
dan konstanta $M > 0$ sedemikian sehingga $\sbsb fT(x) \le M$ untuk setiap $T \in \ofml T$ dan setiap
$x \in U$. Pilih titik $a \in B$ dan bilangan $r > 0$ sedemikian sehingga $\clo{B_r(a)} \subseteq
U$. Kemudian verifikasi bahwa
  \[ \norm{Ty} \le r^{-1}\bigl(\, \sbsb fT(a + ry) + \sbsb fT(a)\,\bigr) \le 2Mr^{-1} \]
untuk setiap $T \in \ofml T$ dan setiap $y \in B$ sedemikian sehingga $\norm y \le 1$. \ns
\end{proof}

\begin{thm}[Banach-Steinhaus] Misalkan $B$ ruang Banach, $W$ ruang linear bernorma,
 \index{teorema Banach-Steinhaus}%
dan $(T_n)$ barisan pemetaan linear terbatas dari $B$ ke~$W$. Jika limit titik demi titik
$\lim_{n \sto \infty} T_nx$ ada untuk setiap $x$ di~$B$, maka pemetaan $S\colon B \sto
W\colon x \mapsto \lim_{n \sto \infty} T_nx$ adalah pemetaan linear terbatas.
\end{thm}

\begin{proof}[\emph{Petunjuk pembuktian}] Untuk setiap $x \in B$, barisan $\bigl(\norm{T_nx}\bigr)_{n=1}^\infty$
konvergen, dan karena itu terbatas. Gunakan~\ref{thm_PUB}.   \ns
\end{proof}

\begin{defn} Suatu subhimpunan $A$ dari ruang linear bernorma $V$ disebut
 \index{terbatas!secara lemah!himpunan dalam ruang linear bernorma}%
 \index{lemah!keterbatasan!dalam ruang linear bernorma}%
 \index{w@$w$-terbatas}%
\df{terbatas secara lemah} ($w$-terbatas) jika $f^{\sto}(A)$ adalah subhimpunan terbatas dari medan skalar
$V$ untuk setiap $f \in V^*$.
\end{defn}

\begin{prop} Suatu subhimpunan $A$ dari ruang Hilbert $H$ terbatas secara lemah jika dan hanya jika untuk setiap $y \in H$
 \index{lemah!keterbatasan!dalam ruang Hilbert}%
himpunan $\{\langle a,y \rangle \colon a \in A\}$ terbatas.
\end{prop}

Proposisi sebelumnya kurang lebih sudah jelas. Proposisi berikutnya merupakan konsekuensi penting dari
\emph{prinsip keterbatasan seragam}.

\begin{prop}\label{C070131} Suatu subhimpunan dari ruang linear bernorma terbatas secara lemah jika dan hanya jika
subhimpunan itu terbatas.
\end{prop}

\begin{proof}[\emph{Petunjuk pembuktian}] Jika $A$ subhimpunan yang terbatas secara lemah dari ruang linear bernorma $V$
dan $f \in V^*$, apa yang dapat Anda katakan tentang $\sup\{\abs{a^{**}(f)}\colon a \in A\}$?  \ns
\end{proof}

\begin{exer} Sahabat baik Anda, Fred R. Dimm, kembali memerlukan bantuan. Kali ini ia mencemaskan
pernyataan (lihat~\ref{C070131}) bahwa suatu subhimpunan ruang linear bernorma terbatas jika
dan hanya jika subhimpunan itu terbatas secara lemah. Ia meninjau barisan $(a^n)$ dalam ruang Hilbert
$l^2$ yang didefinisikan oleh $a^n = \sqrt{n}\,\vc e^n$ untuk setiap $n \in \N$, dengan $\{\vc
e^n\colon n \in \N\}$ basis ortonormal biasa untuk~$l_2$ (lihat
contoh~\ref{C063149}). Ia melihat bahwa barisan itu jelas tidak terbatas (dengan topologi biasa
pada~$l^2$). Namun, baginya barisan itu tampak terbatas secara lemah. Jika tidak, menurut argumennya, maka
berdasarkan \emph{teorema Riesz-Fr\'echet} \ref{000342} akan terdapat barisan $x$ dalam
$l^2$ sedemikian sehingga himpunan $\{\abs{\langle a^n,x \rangle}\colon n \in\N\}$ tidak terbatas. Baginya ini tampak
tidak mungkin terjadi karena barisan semacam itu harus menurun lebih lambat daripada
barisan $\bigl(n^{-1/2}\bigr)$, yang sudah menurun terlalu lambat untuk menjadi anggota~$l^2$.
Tenangkan Fred dengan menemukan barisan yang sesuai. \emph{Petunjuk:} Gunakan barisan yang
memiliki banyak nol.
\end{exer}

\begin{defn} Suatu barisan $(x_n)$ dalam ruang linear bernorma $V$ disebut
 \index{lemah!barisan Cauchy}%
 \index{Cauchy!barisan!lemah}%
 \index{barisan!Cauchy lemah}%
\df{Cauchy secara lemah} jika barisan $(f(x_n))$ bersifat Cauchy untuk setiap $f$ dalam~$V^*$. Barisan
tersebut
 \index{lemah!konvergensi!barisan dalam ruang linear bernorma}%
 \index{konvergensi!lemah}%
 \index{barisan!konvergensi lemah dari suatu}%
\df{konvergen secara lemah} ke titik $a \in V$ jika $f(x_n) \sto f(a)$ untuk setiap $f \in V^*$
(yakni, jika barisan itu konvergen dalam topologi lemah yang diinduksi oleh anggota-anggota~$V^*$). Dalam hal ini
kita tulis
 \index{<arrowtow@$x_n \to^w a$ (konvergensi sekuensial lemah)}%
$x_n \to^w a$. Ruang $V$ disebut
 \index{lengkap!secara lemah}%
 \index{lemah!kelengkapan}%
\df{lengkap sekuensial secara lemah} jika setiap barisan Cauchy secara lemah di $V$ konvergen secara lemah.
\end{defn}

\begin{prop} Dalam ruang linear bernorma, setiap barisan Cauchy secara lemah terbatas.
\end{prop}

\begin{prop} Setiap ruang Banach refleksif lengkap sekuensial secara lemah.
\end{prop}

\begin{exam} Misalkan $f_n(t) = \left\{
                             \begin{array}{ll}
                               1 - nt, & \hbox{jika $0 \le t \le \frac1n$;} \\
                               0, & \hbox{jika $\frac1n < t \le 1$.}
                             \end{array}
                           \right.$ untuk setiap $n \in \N$. Barisan fungsi $(f_n)$ menunjukkan bahwa ruang Banach
$\fml C([0,1])$ tidak lengkap sekuensial secara lemah.
\end{exam}

\begin{cor} Ruang Banach $\fml C([0,1])$ tidak refleksif.
\end{cor}

\begin{prop} Misalkan $\fml E$ basis ortonormal untuk suatu ruang Hilbert~$H$ dan $(x_n)$ suatu
barisan di~$H$. Maka $x_n \to^w \vc 0$ jika dan hanya jika $(x_n)$ terbatas dan $\langle
x_n,e \rangle \sto 0$ untuk setiap $e \in \fml E$.
\end{prop}

\begin{defn} Misalkan $H$ ruang Hilbert. Suatu jaring $(T_\lambda)$ dalam $\ofml B(H)$
 \index{lemah!topologi operator!konvergensi dalam}%
 \index{operator!topologi!lemah}%
 \index{topologi!operator lemah}%
 \index{lemah!konvergensi!operator ruang Hilbert}%
\df{konvergen secara lemah} (atau \df{konvergen dalam topologi operator lemah}) ke $S \in \ofml B(H)$ jika
  \[ \langle T_\lambda x,y \rangle \sto \langle Sx,y \rangle \]
untuk setiap $x$, $y \in H$. Dalam hal ini kita tulis
$T_\lambda~\to^{\textbf{WOT}}~S$.
 \index{<arrowtowot@$T_\lambda~\to^{\textbf{WOT}}~S$}%

Suatu keluarga $\ofml T \subseteq \ofml B(H)$ disebut
 \index{lemah!keterbatasan!keluarga operator ruang Hilbert}%
 \index{terbatas!secara lemah}
 \index{terbatas!dalam topologi operator lemah}%
 \index{lemah!topologi operator!keterbatasan dalam}%
 \index{WOT@\textbf{WOT} (topologi operator lemah)}%
\df{terbatas secara lemah} (atau \df{terbatas dalam topologi operator lemah}) jika untuk setiap $x$, $y \in H$
terdapat konstanta positif $\alpha_{x,y}$ sedemikian sehingga
  \[\abs{\langle Tx,y \rangle} \le \alpha_{x,y}\]
untuk setiap $T \in \ofml T$.
\end{defn}

Sayangnya definisi sebelumnya dapat menjadi sumber kebingungan. Karena $\ofml B(H)$ adalah
ruang linear bernorma, ruang itu sudah mempunyai topologi ``lemah'' (lihat~\ref{C067434}). Lalu, bagaimana menangani dua
topologi berbeda pada $\ofml B(H)$ yang mempunyai nama sama? Sebagian penulis menyebut topologi yang baru saja didefinisikan sebagai
\emph{topologi operator lemah}, sehingga mencegah timbulnya kebingungan. Penulis lain berpendapat bahwa
topologi lemah dalam pengertian ruang linear bernorma begitu jarang berguna dalam teori operator pada ruang Hilbert
sehingga, kecuali diberi penjelas khusus, frasa \emph{topologi lemah} dalam konteks operator ruang Hilbert
seharusnya dipahami sebagai topologi yang baru saja didefinisikan. (Lihat, misalnya, komentar Halmos dalam dua
paragraf sebelum Soal 108 di~\cite{Halmos:1982}.) Demikian pula, kita sering melihat frasa \emph{suatu
himpunan terbatas yang terdiri atas operator-operator ruang Hilbert.} Kecuali diberi penjelas lain, yang dimaksud adalah bahwa himpunan
tersebut \emph{terbatas seragam}. Bagaimanapun juga, berhati-hatilah terhadap frasa-frasa ini ketika membaca.

Berikut ini konsekuensi penting dari \emph{prinsip keterbatasan seragam}~\ref{thm_PUB}.

\begin{prop} Setiap keluarga operator ruang Hilbert yang terbatas dalam topologi operator
lemah adalah terbatas seragam.
\end{prop}

\begin{proof}[\emph{Petunjuk pembuktian}] Misalkan $H$ ruang Hilbert dan $\ofml T$ keluarga operator pada~$H$
yang terbatas terhadap topologi operator lemah. Gunakan proposisi~\ref{C070131} untuk menunjukkan, bagi
$y \in H$ yang tetap, bahwa $\{T^*y\colon T \in \ofml T\}$ adalah subhimpunan terbatas dari~$H$. Gunakan hal ini (dan~\ref{C070131}
sekali lagi) untuk menunjukkan bahwa $\{Tx\colon \text{$T \in \ofml T$, $x \in H$, dan $\norm x \le 1$}\}$ terbatas di~$H$. \ns
\end{proof}

Seperti dinyatakan di atas, banyak penulis akan menyingkat pernyataan proposisi sebelumnya menjadi:
\emph{Setiap keluarga operator ruang Hilbert yang terbatas secara lemah adalah terbatas.}

\begin{defn}
Suatu jaring $(T_\lambda)$ dalam $\ofml B(H)$
 \index{kuat!topologi operator!konvergensi dalam}%
 \index{operator!topologi!kuat}%
 \index{topologi!operator kuat}%
 \index{kuat!konvergensi!operator ruang Hilbert}%
 \index{SOT@\textbf{SOT} (topologi operator kuat)}%
\df{konvergen secara kuat} (atau \df{konvergen dalam topologi operator kuat}) ke $S \in \ofml B(H)$ jika
  \[ T_\lambda x \sto Sx \]
untuk setiap $x \in H$. Dalam hal ini kita tulis $T_\lambda~\to^{\textbf{SOT}}~S$.
 \index{<arrowtosot@$T_\lambda~\to^{\textbf{SOT}}~S$}%

Konvergensi norma biasa dari operator-operator dalam $\ofml B(H)$ sering disebut
 \index{seragam!topologi operator!konvergensi dalam}%
 \index{operator!topologi!seragam}%
 \index{topologi!operator seragam}%
 \index{seragam!konvergensi!operator ruang Hilbert}%
\df{konvergensi seragam} (atau \df{konvergensi dalam topologi operator seragam}). Artinya, kita tulis
 \index{<arrowtouniform@$T_\lambda~\sto~S$}%
$T_\lambda \sto S$ jika $\norm{S - T_\lambda}\sto 0$.
\end{defn}

\begin{prop} Misalkan $H$ ruang Hilbert, $(T_n)$ barisan dalam $\ofml B(H)$,
dan $B \in \ofml B(H)$. Maka implikasi-implikasi berikut berlaku:
   \[ T_n \sto B \quad \implies \quad T_n~\to^{\textbf{SOT}}~B
                                \quad \implies \quad T_n~\to^{\textbf{WOT}}~B. \]
\end{prop}

\begin{prop} Misalkan $(S_n)$ dan $(T_n)$ barisan operator pada suatu ruang Hilbert~$H$. \\
Jika $S_n \to^{\textbf{WOT}} A$ dan $T_n \to^{\textbf{SOT}} B$, maka $S_nT_n \to^{\textbf{WOT}} AB$.
\end{prop}

\begin{proof}[\emph{Petunjuk pembuktian}] Tuliskan $S_nT_n - AB$ sebagai $S_nT_n - S_nB + S_nB - AB$.  \ns
\end{proof}

\begin{exam} Dalam proposisi sebelumnya, hipotesis $T_n \to^{\textbf{SOT}} B$ tidak dapat
diganti dengan $T_n~\to^{\textbf{WOT}}~B$.
\end{exam}

\begin{proof}[\emph{Petunjuk pembuktian}] Tinjau pangkat-pangkat operator geser unilateral.  \ns
\end{proof}






\endinput
